\documentclass[output=paper,colorlinks,citecolor=brown]{langscibook}
\ChapterDOI{10.5281/zenodo.22078591}
\author{Barbara Tomaszewicz-Özakın\orcid{0000-0002-6388-8778}\affiliation{RWTH Aachen University}}
% REPLACE THE ABOVE WITH YOU AND YOUR CO-AUTHORS, YOUR AFFILIATIONS, AND YOUR ORCID NUMBERS
% rules for affiliation: If there's an official English version, use that (find out on the official website of the university); if not, use the original
% orcid does not appear printed; it's metainformation used for later indexing

\SetupAffiliations{mark style=none}
% comment in (i.e. remove the % sign at the beginning of the line) the above line if you are a single author or all authors have the same affiliation

%\lehead{Šimík, Gehrke, Lenertová, Meyer, Szucsich \& Zaleska}
% in case the running head with authors exceeds one line (which is the case in this example document), use one the above method to turn it into a single line; otherwise comment the line above out with % and ignore it

\title{Family of exclusives in Polish}
% REPLACE THE ABOVE WITH YOUR PAPER TITLE
% provide a shorter version of your title in case it does not fit a single line in the running head
% in this form: \title[short title]{full title}

%\abstract{Abstract goes here and should not have more than 150 words.
\abstract{I define the semantics of four exclusive adverbials in Polish -- \textit{tylko}, \textit{zaledwie}, \textit{dopiero}, and \textit{aż} -- as composed of three core components: an assertion of exclusivity, a factive presupposition and a scalar presupposition. By minimally adjusting parameters such as the direction and dimension of the scale and the size of the associated focus, different flavors of exclusivity emerge. These parameter settings give rise to temporal readings (\textit{dopiero}), approximative and counterfactual readings (\textit{zaledwie}), and even additive-like readings (\textit{aż}).

\keywords{\textit{only}, \textit{even}, exclusivity, scalarity, assertion, presupposition, focus association}
}

\IfFileExists{../localcommands.tex}{
   \addbibresource{../localbibliography.bib}
   \usepackage{tabularx,multicol}
\usepackage{url}
\urlstyle{same}

 
\usepackage{langsci-optional}
\usepackage{langsci-lgr}
\usepackage{langsci-gb4e}

% ADDED BY VOLUME EDITORS

% \usepackage{langsci-osl} % macros specific to Open Slavic Linguistics; PLACED DIRECTLY IN localcommands.tex
\usepackage{siunitx}
\usepackage[linguistics]{forest}
% \usepackage{caption} %borik (?)
\usepackage{pifont} %geist

% 08_muellerreichau

% \usepackage{caption}
\usepackage{subcaption}
\usepackage{cleveref}

\usepackage{drs}
\usepackage{diagbox}

% 09_simik

\usepackage{multirow}
\usepackage{tabto}

% 10_stegovec

\usetikzlibrary{arrows,patterns}
\usepackage{listings}
% \usepackage{multirow}

% 01_gehrke-simik

% \usepackage{comment}

% 13_wagiel

% \usepackage{pifont} % for checkmarks (\ding{51}, \ding{55})


\usepackage{langsci-branding}

   % \newcommand*{\orcid}{}


\makeatletter
\let\thetitle\@title
\let\theauthor\@author
\makeatother

\newcommand{\togglepaper}[1][0]{
   \bibliography{../localbibliography}
   \papernote{\scriptsize\normalfont
     \theauthor.
     \titleTemp.
     To appear in:
     E. Di Tor \& Herr Rausgeberin (ed.).
     Booktitle in localcommands.tex.
     Berlin: Language Science Press. [preliminary page numbering]
   }
   \pagenumbering{roman}
   \setcounter{chapter}{#1}
   \addtocounter{chapter}{-1}
}

\newbool{bookcompile}
\booltrue{bookcompile}
\newcommand{\bookorchapter}[2]{\ifbool{bookcompile}{#1}{#2}}

\AtBeginDocument{\nogltOffset} %removes extra spacing before trans line in examples - SPECIFIC TO OPEN SLAVIC LINGUISTICS, PREVIOUSLY AGREED UPON


%%%%%%%%%%%%%%%%%%%%%%%%%%%%%%%%%%%%%%%%%%%%%%%%%%%%%%%%%%%%%%%%%%%%%
%%      File: langsci-osl.sty
%%    Author: Radek Simik and Language Science Press (http://langsci-press.org)
%%      Date: 2019-02-18
%%   Purpose: This file contains macros for the Open Slavic Linguistics at Language Science Press series and is included 
%%            into langscibook.cls
%%  Language: LaTeX
%%   Licence: 
%%%%%%%%%%%%%%%%%%%%%%%%%%%%%%%%%%%%%%%%%%%%%%%%%%%%%%%%%%%%%%%%%%%%%

% FIXING GRAY

\definecolor{lsDOIGray}{cmyk}{0,0,0,0.45}

% PACKAGES REQUIRED

\RequirePackage{relsize}
\RequirePackage{stmaryrd}
\RequirePackage{array}
\RequirePackage{tabularx}

% % KEYWORDS

% \newcommand{\keywords}[1]{\noindent\textbf{Keywords:} {#1}}

% SEMANTICS MACROS - GENERAL FOR ALL PAPERS

\newcommand{\sib}[1]{$\llbracket${#1}$\rrbracket$} % = semantic interpretation brackets; puts text into double square brackets
\newcommand{\stb}[1]{\langle{#1}\rangle} % = semantic type brackets; puts stuff into semantic type brackets; necessary to use in the form $\st{}$
\newcommand{\cnst}[1]{\textsf{\relsize{-1}\MakeUppercase{#1}}} %creates sans-serif small-caps, used for typesetting logical constants which have no other appropriate symbols (such as Greek letters)

% MINUS HSPACE MACRO

\newcommand{\minsp}[1]{#1\hspace{-2.5pt}} % use for non-glossed elements in glossed examples (typically stars, brackets, slashes, ...)

% CITATION MACROS

\newcommand{\citeposst}[1]{\citeauthor{#1}'s (\citeyear{#1})} % produces Chomsky's (1995)
\newcommand{\citeposstpg}[2]{\citeauthor{#1}'s (\citeyear[#2]{#1})} % produces Chomsky's (1995: page)
\newcommand{\citepossalt}[1]{\citeauthor{#1}'s \citeyear{#1}} % produces Chomsky's 1995
\newcommand{\citepossaltpg}[2]{\citeauthor{#1}'s \citeyear[#2]{#1}} % produces Chomsky's 1995: page

% BARPLOT

\usepackage{pgfplots}
\pgfplotsset{compat=1.18} 
\newcommand{\barplot}[5]{%
  \begin{tikzpicture}
    \begin{axis}[
	xlabel={#1},  
	ylabel={#2}, 
	axis lines*=left, 
        width  = {#3\textwidth},
	height = .3\textheight,
    	nodes near coords, 
	xtick=data,
	x tick label style={},  
	ymin=0,
	symbolic x coords={#4},
	]
	\addplot+[ybar,lsRichGreen!80!black,fill=lsRichGreen] plot coordinates {
	    #5
	}; 
    \end{axis} 
  \end{tikzpicture} 
}



%%%%%%%%%%%%%%%%%%%%%%%%%%%
%%% 04_filip %%%

% \newcommand{\mC}[1]{\multicolumn{1}{c}{#1}} % multicolumn with a single column; ALREADY DEFINED

\usetikzlibrary{arrows, arrows.meta}
\usetikzlibrary{positioning, decorations, 
                decorations.pathreplacing, calligraphy}

%%%%%%%%%%%%%%%%%%%%%%%%%
%%% 08_muellerreichau %%%

\captionsetup[subfigure]{subrefformat=simple,labelformat=simple}
\renewcommand\thesubfigure{(\alph{subfigure})}

%%%%%%%%%%%%%%%%%%%%%%%%%%%
%%% 09_simik

\newcommand{\citepossts}[1]{\citeauthor{#1}' (\citeyear{#1})} % produces Rawlins' (2008)
\newcommand{\citepossalts}[1]{\citeauthor{#1}' \citeyear{#1}} % produces Rawlins' 2008


%%%%%%%%%%%%%%%%%%%%%%%%%%%%%%%
%%% 10_stegovec

\newcommand{\poscite}[1]{\citeauthor{#1}'s (\citeyear{#1})}

%extra commands
\newcommand{\fst}{\textsc{1p}}
\newcommand{\snd}{\textsc{2p}}
\newcommand{\trd}{\textsc{3p}}
\newcommand{\refl}{\textsc{refl}}
\newcommand{\excl}{\textsc{excl}}
\newcommand{\incl}{\textsc{incl}}
\newcommand{\sg}{\textsc{sg}}
\newcommand{\pl}{\textsc{pl}}
\newcommand{\dl}{\textsc{dl}}
\newcommand{\fem}{\textsc{f}}
\newcommand{\masc}{\textsc{m}}
\newcommand{\neut}{\textsc{n}}
\newcommand{\ImpCl}{IMPERATIVE}


\newcommand{\tikzmark}[1]{\tikz[overlay,remember picture] \node (#1) {};}
\forestset{M/.style={
		for tree={s sep=5mm, inner sep=0.5pt, l=5mm,
			calign=fixed edge angles,
			calign primary angle=-65,
			calign secondary angle=65},
	},
	Mid/.style={
		for tree={s sep=0mm, inner sep=0pt, l=0mm,
			calign=fixed edge angles,
			calign primary angle=-62.5,
			calign secondary angle=62.5},
	},
	word/.style={before typesetting nodes={content=\emph{##1}},
		no edge,l=0,for parent={l sep=0}},
	feat/.style={before typesetting nodes={content={{[}##1{]}}},
		no edge,l=0,for parent={l sep=0}},
	feat2/.style={before typesetting nodes={content={##1}},
		no edge,l=0,for parent={l sep=0}},
	llap/.style={
		tikz+={%
			\edef\forest@temp{\noexpand\node[\option{node options},
				anchor=base east,at=(.base east)]}%
			\forest@temp{#1\phantom{\option{environment}}};
		}
	},
	rlap/.style={
		tikz+={%
			\edef\forest@temp{\noexpand\node[\option{node options},
				anchor=base west,at=(.base west)]}%
			\forest@temp{\phantom{\option{environment}}#1};
		}
	}
}

\makeatletter
\newcommand{\coloruline}[2]{%
	\newcommand\temp@reduline{\bgroup\markoverwith
		{\textcolor{#1}{\rule[-0.5ex]{2pt}{0.4pt}}}\ULon}%
	\temp@reduline{#2}%
}

\newcommand{\coloruuline}[2]{%
	\UL@protected\def\temp@uuline{\leavevmode \bgroup
		\UL@setULdepth
		\ifx\UL@on\UL@onin \advance\ULdepth2.8\p@\fi
		\markoverwith{\textcolor{#1}{\lower\ULdepth\hbox
				{\kern-.03em\vbox{\hrule width.2em\kern1\p@\hrule}\kern-.03em}}}%
		\ULon}%
	\temp@uuline{#2}%
}

\newcommand{\coloruwave}[2]{%
	\UL@protected\def\temp@uwave{\leavevmode \bgroup 
		\ifdim \ULdepth=\maxdimen \ULdepth 3.5\p@
		\else \advance\ULdepth2\p@ 
		\fi \markoverwith{\textcolor{#1}{\lower\ULdepth\hbox{\sixly \char58}}}\ULon}
	\font\sixly=lasy6 % does not re-load if already loaded, so no memory drain.
	\temp@uwave{#2}%
}
\makeatother

%%%%%%%%%%%%%%%%%
%%% 13_wagiel %%%

\newcommand\irregularcircle[2]{% radius, irregularity
  \pgfextra {\pgfmathsetmacro\len{(#1)+rand*(#2)}}
  +(0:\len pt)
  \foreach \a in {10,20,...,350}{
    \pgfextra {\pgfmathsetmacro\len{(#1)+rand*(#2)}}
    -- +(\a:\len pt)
  } -- cycle
}


% IF POSSIBLE, CAN WE USE THIS SETTING JUST FOR 13_wagiel?

% \forestset{
% default preamble={for tree={s sep=10mm, inner sep=0, l=0}},
% fairly nice empty nodes/.style={
%         delay={where content={}{shape=coordinate,
%               for siblings={anchor=north}}{}},
% for tree={s sep=4mm}},
% }

   %% hyphenation points for line breaks
%% Normally, automatic hyphenation in LaTeX is very good
%% If a word is mis-hyphenated, add it to this file
%%
%% add information to TeX file before \begin{document} with:
%% %% hyphenation points for line breaks
%% Normally, automatic hyphenation in LaTeX is very good
%% If a word is mis-hyphenated, add it to this file
%%
%% add information to TeX file before \begin{document} with:
%% %% hyphenation points for line breaks
%% Normally, automatic hyphenation in LaTeX is very good
%% If a word is mis-hyphenated, add it to this file
%%
%% add information to TeX file before \begin{document} with:
%% \include{localhyphenation}
\hyphenation{
    par-a-digm
    Cart-wright %filip
    Truswell %docekal
    Shcher-bi-na %simik
    Mün-chen %gehrke
    Mo-ni-ka %gehrke
    spi-sov-né %gehrke
    li-te-ra-tur-no-go %gehrke
    russ-koj %gehrke
    so-ver-šen-no-go %gehrke
    russ-kom %gehrke
    sla-vjan-sko-mu %gehrke
    Zeit-schrift %gehrke
    Na-ta-sha %gehrke
    Zu-stands-pas-siv %gehrke
    Um-gangs-spra-che %gehrke
    Bra-so-vea-nu %geist
}

\hyphenation{
    par-a-digm
    Cart-wright %filip
    Truswell %docekal
    Shcher-bi-na %simik
    Mün-chen %gehrke
    Mo-ni-ka %gehrke
    spi-sov-né %gehrke
    li-te-ra-tur-no-go %gehrke
    russ-koj %gehrke
    so-ver-šen-no-go %gehrke
    russ-kom %gehrke
    sla-vjan-sko-mu %gehrke
    Zeit-schrift %gehrke
    Na-ta-sha %gehrke
    Zu-stands-pas-siv %gehrke
    Um-gangs-spra-che %gehrke
    Bra-so-vea-nu %geist
}

\hyphenation{
    par-a-digm
    Cart-wright %filip
    Truswell %docekal
    Shcher-bi-na %simik
    Mün-chen %gehrke
    Mo-ni-ka %gehrke
    spi-sov-né %gehrke
    li-te-ra-tur-no-go %gehrke
    russ-koj %gehrke
    so-ver-šen-no-go %gehrke
    russ-kom %gehrke
    sla-vjan-sko-mu %gehrke
    Zeit-schrift %gehrke
    Na-ta-sha %gehrke
    Zu-stands-pas-siv %gehrke
    Um-gangs-spra-che %gehrke
    Bra-so-vea-nu %geist
}

   \boolfalse{bookcompile}
   \togglepaper[12]%%chapternumber
}{}

\begin{document}
\maketitle

%\input{template-osl.tex}
% Just comment out the line above (place % at its beginning) and start writing your own document
% WRITE DIRECTLY HERE, INTO main.tex, DO NOT START YOUR OWN tex FILE!

% START WRITING HERE

\section{Introduction}

This chapter presents the behavior of four exclusive adverbials in Polish, \textit{tylko}, \textit{zaledwie}, \textit{dopiero} and \textit{aż}, and shows how they can be defined as a family: Each lexical entry includes a common core (assertion of exclusivity, presupposition of scalarity) with minimal differences in specific parameters (the direction and the dimension of the scale, the size of the associated focus).\footnote{The glosses of the Polish particles in the examples reflect the parallels with the English counterparts. In \sectref{sec:tom:4}, we will see that \textit{tylko} is similar to \textit{only} and \textit{zaledwie} to \textit{merely}. \textit{Dopiero} has a German counterpart, \textit{erst}, which has been named \textsc{temporal} \textit{only}. \textit{Aż} on my view is a scalar opposite of \textit{merely}, so it has no approximate translation and is glossed here as \textsc{aż}.} We will start with a context where all four adverbs have a similar effect on the meaning and then use a battery of tests to identify the individual meaning components.

All of the adverbs contribute to the sentence in \REF{ex:tom:1} the additional meaning that the 30\textsuperscript{th} place is a weak result, i.e. the team did not perform any better than that.

\ea\label{ex:tom:1} 
\gll Zespół był     słaby  i         zajął         \minsp{\{} tylko / zaledwie / dopiero / aż\}   30   miejsce.\\
team    was  poor    and    achieved  {} only {} merely {} only\textsubscript{\textsc{temp}} {} \textsc{aż}     30\textsuperscript{th}  place\\
\glt `The team performed poorly and earned only 30\textsuperscript{th} place.'
\z

\noindent %The 
English \textit{only} has exactly this meaning `(at least) $X$ and nothing more', with `nothing more' excluding all contextually ranked higher alternatives. Is it only a coincidence that in this particular context the four adverbials seem to mean `only'? Could the sentence rather be saying that ‘The team was ranked as low as the 30\textsuperscript{th} place’? This meaning is similar but it does not contain the exclusive component `nothing more'. I argue that the four adverbial particles should all be classified as exclusives, i.e., they all share the exclusive meaning component, and that this primary component is at-issue (\citealt{Potts2005}, \citealt{beaverclark08}, \citealt{Simons.etal2011}, \citealt{roberts11}). Each one of them, however, has its own flavor so that in most contexts they are not interchangeable as they are in \REF{ex:tom:1}. I identify the differences in their use that make them similar only under certain conditions. It turns out that Polish does not have one adverbial that can contribute the exclusive meaning in all contexts, the way %the 
English \textit{only} can. Instead, each of the items places particular restrictions on the dimension of scale of the excluded alternatives. Furthermore, the direction of exclusion differs: alternatives both higher and lower on the scale can be excluded \citep{tomaszewicz12,tomaszewicz13,tomaszewicz13a}. I propose lexical entries for each exclusive such that they all include a common core with minimal differences in specific parameters.

\section{Exclusivity and scalarity}
\label{sec:tom:2}

The adverb \textit{only} contributes the meaning of exclusion: in contrast to \REF{ex:tom:2}, the sentence in \REF{ex:tom:3} states that Maria met Jan and she did not meet anyone else. 

\ea\label{ex:tom:2}
Maria met Jan.
\z 

\ea\label{ex:tom:3} 
Maria met only Jan.
\z

\noindent In some contexts \textit{only} additionally contributes scalarity by implying a ranking among the excluded alternatives. Consider how  in \REF{ex:tom:4} we understand that the 30th place is not a good result, even though we do not know how many places were assigned in total. In \REF{ex:tom:5} it is implied that 9 PM is late, and in \REF{ex:tom:6}--\REF{ex:tom:7} that the manager is lower on the scale than some contextually relevant alternative.

\ea\label{ex:tom:4}
Jan got only the 30th place.
\z

\ea\label{ex:tom:5}
Jan arrived only at 9 PM.
\z

\ea\label{ex:tom:6}
Jan is only/merely a manager.
\z

\ea\label{ex:tom:7}
Maria talked only/merely to the manager.
\z

\noindent In the example in \REF{ex:tom:4} it is obvious that you can be ranked at a single place only, so the exclusive meaning should be paraphrased as `Jan got no better place than the 30th' rather than `Jan got the 30th place and no other place' in contrast to the paraphrase of the sentence in \REF{ex:tom:3}. These two readings of \textit{only} have been called \textsc{scalar} and \textsc{quantificational} (\citealt{bonomicasalegno93}, \citealt{krifka93}, \citealt{klinedinst05}, \citealt{riester06}, \citealt{beaverclark08}, \citealt{zeevat09}, among others). The scalar reading of the sentence in \REF{ex:tom:5} is roughly equivalent to `Jan arrived no sooner than 9 PM' because the most plausible scenario is that Jan arrived once, as opposed to `Jan arrived at 9 PM and at no other time', which would be the quantificational reading. In contrast to \REF{ex:tom:4}, in \REF{ex:tom:6} Jan is a manager in addition to being other obvious things such as human, terrestrial, mortal… These properties are not ordered on the scale with respect to the property of being a manager, they are not contextually relevant. Thus, in \REF{ex:tom:6}, \textit{only} evokes a contextually relevant scale and places ‘manager’ low on that scale, as illustrated in \REF{ex:tom:8}. The example in \REF{ex:tom:7} also shows that the scale of exclusion involves only relevant alternatives, i.e., the people ranked in terms of importance in \REF{ex:tom:8}, because \REF{ex:tom:7} is true even if Maria talked to the manager’s secretary (assuming that this conversation was irrelevant to the topic of her conversation with the manager). Note that in English \textit{merely} is dedicated to scalar readings, \REF{ex:tom:6}--\REF{ex:tom:7} \citep{coppockbeaver11, coppockbeaver14}.

\ea\label{ex:tom:8}
Scale for \REF{ex:tom:6}--\REF{ex:tom:7}:\\
`Jan is the president.'\\
`Jan is a vice president.'\\
...\\
`Jan is a senior manager.'\\
`Jan is a manager.'\\
`Jan is an assistant manager.'\\
...\\
`Jan is an employee.'
\z 

\noindent The Polish exclusives in the example in \REF{ex:tom:1}, \textit{tylko}, \textit{zaledwie}, \textit{dopiero} and \textit{aż}, are all compatible with the scalar readings, but now the examples in \REF{ex:tom:9}--\REF{ex:tom:12} will show that Polish does not have one adverbial that can contribute the exclusive meaning in all contexts like %the 
English \textit{only} can.\footnote{All non-English examples in this chapter are from Polish, unless indicated otherwise.} The quantificational reading requires \textit{tylko}, as shown in the sentence in \REF{ex:tom:9} which is the counterpart of the sentence we saw in \REF{ex:tom:3}. With temporal scales, on the other hand, \textit{tylko} cannot be used, \REF{ex:tom:10}, unlike \textit{only} in \REF{ex:tom:5}. With the scale of important people, \REF{ex:tom:11}--\REF{ex:tom:12}, \textit{tylko}, \textit{zaledwie} and \textit{dopiero} are possible, though here the manager is low on the scale, just like in \REF{ex:tom:6}--\REF{ex:tom:7} with \textit{only/merely}.

\ea\label{ex:tom:9}
\gll Maria   spotkała \minsp{\{} tylko / \minsp{\#} zaledwie / \minsp{\#} dopiero / \minsp{\#} aż\}   Jana.\\
Maria  met {} only {} {} merely {} {}  only\textsubscript{\textsc{temp}} {} {} \textsc{aż}            Jan\\
\glt ‘Maria met only Jan and nobody else.’
\z

\ea\label{ex:tom:10}
\gll Jan   przyjechał  \minsp{\{} dopiero / aż / \minsp{\#} tylko / \minsp{?} zaledwie\}   o   9-tej.\\
Jan arrived {} only\textsubscript{\textsc{temp}} {} \textsc{aż} {} {} only {} {} merely at 9\\
\glt ‘Jan arrived only at 9 PM and no sooner.’
\z 

\ea\label{ex:tom:11}
\gll Jan jest \minsp{\{} tylko / zaledwie / dopiero /  \minsp{\#} aż\}     menedżerem.\\
Jan is {} only {} merely {}  only\textsubscript{\textsc{temp}} {} {} \textsc{aż} manager\\
\glt ‘Jan is only/merely a manager (so far) and nobody more important.’
\z 

\ea\label{ex:tom:12}
\gll Maria rozmawiała \minsp{\{} tylko / zaledwie / dopiero / \minsp{\#} aż\} z menedżerem.\\
Maria talked {} only {} merely {} only\textsubscript{\textsc{temp}} {} {} \textsc{aż} with manager\\
\glt ‘Maria only/merely talked to the manager (so far) and nobody more important.’
\z 

\noindent Crucially, the scalar readings of \REF{ex:tom:11}--\REF{ex:tom:12} “flip” when \textit{aż}  is used: In \REF{ex:tom:13} the manager is interpreted to be high on the scale of importance. The exclusive \textit{aż} results in a reading that is similar to the effect of the addition of \textit{even} (Polish \textit{nawet}). In the sentence with \textit{nawet} in \REF{ex:tom:14}, similarly to \REF{ex:tom:13}, the manager is implied to be a noteworthy person for Maria to talk to. However, \textit{nawet} is not compatible with the continuation that Maria did not talk to anybody else, \REF{ex:tom:14}, in contrast to \textit{aż} in \REF{ex:tom:15}. This is because \textit{nawet}/\textit{even} is an additive scalar particle (\citealt{horn69}, \citealt{karttunenpeters79}, \citealt{rooth85,rooth92}, a.o.), whereas \textit{aż} is an exclusive as I argued in \citet{tomaszewicz12,tomaszewicz13,tomaszewicz13a}.

\ea\label{ex:tom:13}
\gll Maria   rozmawiała   \minsp{\{\#} tylko / \minsp{\#} zaledwie / \minsp{\#} dopiero / aż\}   z       menedżerem.\\
Maria talked {} only {} {} merely {} {} only\textsubscript{\textsc{temp}} {} \textsc{aż} with  manager\\
\glt ‘Maria talked to somebody so important as the manager.’
\z 

\ea\label{ex:tom:14}
\gll Maria rozmawiała nawet z menedżerem, \minsp{\#} ale nie rozmawiała z nikim innym.\\
Maria talked even with  manager {} but not talked with nobody else\\
\glt `Maria even talked to the manager, \textit{[Incongruous:]} but she did not talk to anybody else.’
\z 

\ea\label{ex:tom:15}
\gll Maria rozmawiała aż z menedżerem, ale nie rozmawiała z nikim innym. \\
Maria talked \textsc{aż} with  manager but not talked with  nobody else\\
\glt `Maria talked to somebody so important as the manager, but she did not talk to anybody else.’
\z 


\noindent I proposed in \citet{tomaszewicz12,tomaszewicz13,tomaszewicz13a} that \textit{aż} is the scalar opposite of \textit{merely}/\textit{zaledwie}: It asserts exclusion just like \textit{only}/\textit{merely}, evokes scales like \textit{merely}, but it places the mentioned alternative high on the scale. In this chapter, I provide tests to specify the meaning components for each of the exclusives shown in \REF{ex:tom:1}, demonstrating that the semantics contains a common core predicting the existence of scalar opposites like \textit{aż} cross-linguistically, even though such a lexical item is unavailable in, for instance, English.

Our starting point is the basic meaning contribution of \textit{only}/\textit{merely}, which involves, on the one hand, exclusivity \REF{ex:tom:16a}--\REF{ex:tom:16b} and, on the other hand, the prejacent, which is the meaning of the sentence without \textit{only}/\textit{merely} \REF{ex:tom:16c} (\citealt{klinedinst05}, \citealt{roberts06}, \citealt{vanrooijschulz07}, \citealt{beaverclark08}, a.o.).\footnote{This use of the term \textsc{prejacent} was introduced by \citet{horn96}; however, \citet{roberts06} observes that often with \textit{only} the prejacent does not correspond to any overt syntactic constituent. \citet{vonfintelheim02} use the term for the constituent that is the sister of the operator at the level of logical syntax. \citet{roberts06} addresses the debate on the status of the prejacent with \textit{only} arguing that it is “closer to a presupposition than to an entailment or a conversational implicature”. \citet{crnic12} argues that the mixed behavior (prejacent presupposition vs. existential presupposition) results from two separate meaning components at LF with different scope.}

\ea\label{ex:tom:16}
Maria only talked to the manager.\smallskip\\
\textit{Exclusive/negative contribution}:
\ea\label{ex:tom:16a}   
`Maria did not talk to anybody else than the manager.'\\\hfill \textsc{quantificational reading}
\ex\label{ex:tom:16b}  
`Maria did not talk to anybody more noteworthy than a manager.'\\\hfill \textsc{scalar reading}
\z
\textit{Positive contribution}:
\begin{xlist}
\exi{c.}{`Maria talked to (at least) the manager.'}\hfill \textsc{prejacent}\label{ex:tom:16c}

\end{xlist} 
\z 

\noindent The quantificational reading in \REF{ex:tom:16a} can be seen as an instance of exclusion along a logical/entailment scale (\citealt{bonomicasalegno93}, \citealt{vanrooij02}, \citealt{beaver04}, \citealt{klinedinst05}, \citealt{riester06}, a.o.). The dimension of the scale is set relevant to the context, and where no such pragmatic ordering is readily available (such as the ranking of people in terms of importance), the logical scale is the default (e.g., all other people are excluded).

In the following sections I present a battery of tests to demonstrate that with all four Polish adverbial particles, \textit{tylko}, \textit{zaledwie}, \textit{dopiero} and \textit{aż}, exclusivity is asserted, the prejacent is presupposed, and scalarity, i.e., the `low/high on the scale' reading, is presupposed. In addition to the exclusive assertion and the prejacent presupposition, it will be shown that scalarity comes from focus association,  so that the focal presupposition supplies the set of alternatives for the ranking.

\section{Focus association}
\label{sec:tom:3}

\textit{Only} belongs to so-called \textsc{focus-sensitive} adverbs because its surface position interacts with the placement of focus. The semantic effects of focus on the interpretation of \textit{only} go hand in hand with syntactic/phonological effects, i.e., the way \textit{only} modifies the meaning of the sentence depends on which constituent is focused, although the syntactic associate needs not always be marked by easily perceptible prosodic prominence. In Polish focus is marked prosodically, with pitch and intensity as the main acoustic correlates \citep{hamlaoui+19}. The focused constituent is called the \textsc{syntactic associator}. It can be a constituent that is immediately modified by \textit{only}, as we have seen above and as shown below in \REF{ex:tom:17}: `only for a loan'. 

\ea\label{ex:tom:17}
\gll Hanka   poprosiła   tylko [\textsubscript{PP} o       [\textsubscript{DP} pożyczkę]],   i       o     nic           więcej / innego.\\  
Hanka asked only {} for {} loan and for nothing more {} else  \\
\glt `Hanka only asked for a loan and nothing more/else.’
\z 

\noindent But the associator can also be more distant, as can be seen in \REF{ex:tom:18}. The different continuations in \REF{ex:tom:18a}--\REF{ex:tom:18d} result from different focus association configurations (associated with different prosodic patterns). Here \textit{tylko} precedes the verb but it can associate with the object DP so that the sentence has the exact same interpretation as in \REF{ex:tom:17}. It can also associate with the verb itself, as shown by the compatibility with the continuation in \REF{ex:tom:18b}, and it can associate with the whole VP, as evidenced by the continuation in \REF{ex:tom:18c}. In contrast to English, \textit{tylko} in the preverbal position can also associate with the subject DP, \REF{ex:tom:18d}. When \textit{tylko} precedes the subject it can, unsurprisingly, associate with it, \REF{ex:tom:19a}, but it can also associate with the whole sentence, \REF{ex:tom:19b}, which, again, is not possible in English.

\ea\label{ex:tom:18}
\gll {}[\textsubscript{DP} Hanka]   tylko [\textsubscript{VP} [\textsubscript{V} poprosiła]   o       [\textsubscript{DP} pożyczkę]],    \\
{} Hanka only {} {} asked for {} loan \\  
\ea\label{ex:tom:18a} 
\gll  … i o nic innego nie poprosiła.\\    
{} and for nothing else not asked\\\hfill \textsc{object} \textsc{DP} \textsc{focus} 
\glt `Hanka asked only for a loan and did not ask for anything else.’ \\
\ex\label{ex:tom:18b}
\gll … a nie zażądała.\\
{} and not demanded\\\hfill \textsc{V} \textsc{focus}
\glt `Hanka only asked for a loan and did not demand it.’
\ex\label{ex:tom:18c}
\gll … i nic innego nie zrobiła.\\
{} and nothing else not did\\\hfill \textsc{VP} \textsc{focus}
\glt `Hanka only asked for a loan and did not do anything else.’
\ex\label{ex:tom:18d}      
\gll … i nikt inny nie prosił.\\
 {} and nobody else not asked\\\hfill \textsc{subject DP focus}
\glt `Only Hanka asked for a loan and nobody else did.’
\z 
\z 

\ea\label{ex:tom:19}
\gll \minsp{(} To) tylko [\textsubscript{S} [\textsubscript{DP} Hanka] poprosiła o pożyczkę],\\    
{} \textsc{top} only {} {} Hanka asked for loan\\      
\ea\label{ex:tom:19a}
\gll … i nikt inny nie prosił.\\
{} and nobody else not asked\\\hfill \textsc{subject} \textsc{DP} \textsc{focus}
\glt `Hanka only asked for a loan and nobody else did.’ 
\ex\label{ex:tom:19b} 
\gll … i nic innego się nie stało.\\                
{} and nothing else self not happened\\\hfill \textsc{sentence} \textsc{focus}
\glt `It is only that Hanka asked for a loan and nothing else happened.’
\z 
\z 

\noindent The associate determines the implicit comparison with alternatives of the same type. So when \textit{tylko} in the sentence-initial position associates with the whole clause (sentence focus in \REF{ex:tom:19b}), what is excluded are other things that could have happened that are ranked higher on the scale of importance. A typical context for \REF{ex:tom:19b} is a question ‘Did anything happen?’ and an answer such as ‘No worries. It is just …’. 

In \REF{ex:tom:20} I provide an example from a Polish comedy TV series ``Alternatywy 4'' (episode 9), from a scene where what appears to be a woman but is in fact a robot, jumps down from a balcony, breaks her neck severing the head, then stands up and picks up the head. The people witnessing this are screaming in shock and a man is trying to calm them down saying that nothing out of the ordinary is happening.

\ea\label{ex:tom:20}
\gll To nic       takiego, to nic         takiego! To tylko ten psychiczny dostał szału!\\
 it   nothing such       it nothing such       it   only     this lunatic       got     amok\\
\glt `It is nothing, it is nothing! It is just that lunatic going crazy!’
\z 

\noindent I am pointing out the differences between Polish and English because the syntax (and prosody) of focus association may vary considerably between languages. The extra options available in Polish are not limited to Slavic languages but are also possible in German (\citealt{bueringhartmann01}, \citealt{steindl13,steindl17}). In \REF{ex:tom:21} German \textit{nur} `only' associates with the subject DP that precedes it, while in \REF{ex:tom:22} \textit{nur} is in the sentence-initial position and can associate either with the subject or the whole clause. In English, on the other hand, the associates of \textit{only} must be in its surface c-command domain (\citealt{taglicht84}, \citealt{beaverclark08}, \citealt{toosarvandani10}, a.o.).

\ea\label{ex:tom:21}
\gll {}[\textsubscript{DP} {Juliane}]\textsubscript{F}   hat nur     die Blumen gegossen.\\                                           {} Juliane    has only    the flowers watered \\   
\glt `It was only Juliane that watered the flowers.’\hfill \citep[German;][]{steindl17}
\z 

\ea\label{ex:tom:22}
\gll Nur [\textsubscript{S} [\textsubscript{DP} {Juliane}]\textsubscript{F$_a$} spielt Klavier]\textsubscript{F$_b$}.\\                           
only {} {} Juliane plays piano\\
\ea `Only Juliane is playing the piano.’ \hfill \textsc{subject DP focus}
\ex `Nothing is going on except for Juliane playing the piano.’ \\\hfill \textsc{sentence focus} 
\z \hfill\citep[German;][]{steindl13}
\z 

\noindent Of course, the different focus association patterns are associated with different prosodic contours of the sentence; specifically, the location of pitch accent and nuclear stress identify the constituents in focus. Without going into details of phonology, syntactic evidence for obligatory focus association comes from clitics. Clitics and reduced pronouns force a wider domain reading (\REF{ex:tom:23b} vs. \REF{ex:tom:23a}), which shows that \textit{tylko} cannot associate with weak (unfocused) pronouns. In \REF{ex:tom:23b} \textit{tylko} has to associate with the adverbial or the verb, which can be independently focused (cf. \citealt{hoeksemazwarts91}, \citealt{beaverclark08} on focus association of \textit{only} in Dutch).

\ea\label{ex:tom:23}
\ea\label{ex:tom:23a}  
\gll {}[\textsubscript{Adv} Wieczorem] tylko widziałem [\textsubscript{DP} niego].\\ 
{} evening only saw {} him\\  
\ea `In the evening I saw him and nobody else.’
\ex `In the evening, and at no other time, I saw him.’
\z 
\ex\label{ex:tom:23b}  
\gll {}[\textsubscript{Adv} Wieczorem]   tylko [\textsubscript{VP} go [\textsubscript{V} widziałem]].\\
{} evening  only {} him {} saw\\
\ea `In the evening I saw him and did not e.g. hear him.’
\ex `In the evening, and at no other time, I saw him.’
\ex \textit{Unavailable:} ‘In the evening I saw him and nobody else.’
\z 
\z 
\z 

\noindent The syntactic tests above, the interaction between the size of the focused constituent and the syntactic position of \textit{tylko} \REF{ex:tom:18}--\REF{ex:tom:19}, and the inability of \textit{tylko} to associate with non-focused elements such as clitics \REF{ex:tom:23}, suggest that \textit{tylko} obligatorily associates with focus.~The same can be shown for other Polish exclusives, as demonstrated for \textit{aż} in \citet{tomaszewicz12,tomaszewicz13,tomaszewicz13a}.

\section{Different ways to exclude alternatives}
\label{sec:tom:4}

\subsection{\textit{tylko}}
\label{sec:tom:4.1}

\textit{Only} is standardly assumed to contribute both to the assertion and the presupposition of the sentence in which it occurs (\citealt{horn69}, \citealt{klinedinst05}, \citealt{roberts06,roberts11}, \citealt{vanrooijschulz07}, \citealt{beaverclark08}). The asserted meaning can be either scalar (the dedicated use of \textit{merely}) or quantificational. In \REF{ex:tom:11}--\REF{ex:tom:12} we saw that \textit{tylko} like \textit{only} can have both scalar and quantificational readings; \REF{ex:tom:12} is repeated below in \REF{ex:tom:24}, and the examples for the two types of scales are given in \REF{ex:tom:25}. In \REF{ex:tom:24} \textit{tylko} associates with the focus on the PP `with the manager'.

\ea\label{ex:tom:24}
\gll Maria rozmawiała tylko z menedżerem.\\
Maria talked only with manager\\
\glt `Maria only talked to the/a manager.’\\
%\textrm{⤳} 
\textsc{Scalar} \textsc{reading}: `Maria talked to nobody ranked higher than the/a manager.'\\
%\textrm{⤳} 
\textsc{quantificational} \textsc{reading}: `Maria talked to nobody other than the/a manager.'
\z 

%\ea\label{ex:tom:25}
%Possible scales for \REF{ex:tom:24}:
%\z

\ea\label{ex:tom:25}
Possible scales for \REF{ex:tom:24}:
\begin{tabularx}{.96\textwidth}{l@{~~}ll@{~~}X}
a.& Contextual&b.& Logical\\
&{`Maria talked with the president'}&&{`Maria talked with 10 people'}\\
&{`Maria talked with vice president'}&&{`Maria talked with 9 people'}\\
&{ …}&&{ …}\\
&{`Maria talked with a manager'}&&{`Maria talked with 2 people'}\\
&{\dots}&&{`Maria talked with 1 person'}\\

&{`Maria talked with an employee'}&&`Maria talked with nobody'\\
\end{tabularx}
\z

% \begin{table}[]
%     \centering
%     \begin{tabular}{ll}
%     \lsptoprule
%          (a) Contextual & (b) Logical \\
%     \midrule
% Maria talked with the president & Maria talked with 10 people\\
% Maria talked with vice president &Maria talked with 9 people\\
% … & …\\
% Maria talked with a senior manager & Maria talked with 2 people\\
% \underline{Maria talked with a manager} & \underline{Maria talked with 1 person}\\
% Maria talked with an assistant manager &Maria talked with nobody\\
% … &\\
% Maria talked with an employee &\\
% \lspbottomrule
%     \end{tabular}
%     \caption{Possible scales for \REF{ex:tom:24}}
%     \label{tab:tom:25}
% \end{table}

\noindent Quantificational \textit{only} asserts that no other alternative is true, e.g., Maria only talked to the manager and nobody else. As noted in the previous section, the alternatives are either seen as not ordered or as an instance of exclusion along a logical/entailment scale, such as a numeral scale, as shown in (\ref{ex:tom:25}b) (\citealt{vanrooij02}, \citealt{klinedinst05}, \citealt{riester06}, \citealt{coppockbeaver11}, a.o.). The dimension of the scale is set pragmatically, and in contexts where no pragmatic ordering is available, the logical scale is the default.

Using tests for presupposition, such as embedding under a question operator or negation (\citealt{chierchiamcconnellginet90}), we can identify three meaning components of \textit{only}, one asserted (exclusivity) and two presupposed components (the prejacent and the ‘low on the scale’ reading).\footnote{Note that there is extensive debate on the nature of the non-truth conditional, i.e., projective content. Many different subtypes have been discussed, including pragmatic and semantic presuppositions \citep{stalnaker73}, conversational and conventional implicature (\citealt{grice81}, \citealt{bach94}, \citealt{Potts2005,potts15}), F-implicature \citep{horn13}, proffered content \citep{roberts96} and several other types and their groupings (e.g., Classes A--D in \citealt{tonhauser+13}).} 
We will use these tests on the three meaning components of the sentence in \REF{ex:tom:24} identified in (\ref{ex:tom:26}a--c). The asserted component is the exclusion of all alternatives higher on the scale: the rank of the person that Maria talked to is no higher than a manager, \REF{ex:tom:26a}.\footnote{The alternative `Maria talked to nobody' is also excluded, which, as suggested by Radek Šimík, could be seen as lower, and not higher, on the scale. However, the logical/entailment scale on the quantificational reading is not one-dimensional but it is a mereological structure (\citealt{riester06}, \citealt{coppockbeaver11}).} \noindent Additionally, \textit{only/merely} contributes two presuppositions, \REF{ex:tom:26b}, \REF{ex:tom:26c}. First, it presupposes that the alternatives under consideration are at least as high on the scale as the prejacent (\citealt{klinedinst05}, \citealt{beaverclark08}). Other relevant ranks of the person that Maria talked to are at least as high as the manager, \REF{ex:tom:26b} (lower ranks are simply not relevant). The second presupposition is that the prejacent is low on the scale of alternative propositions, which I paraphrase as ‘is not significant’, \REF{ex:tom:26c} \citep{klinedinst05}. 


\ea\label{ex:tom:26}
\gll Maria rozmawiała tylko z menedżerem.\\
 Maria talked only with manager\\
\glt `Maria only talked to the manager.’
%\textrm{⤳} 
\ea\label{ex:tom:26a}
Maria talked to nobody ranked higher than a manager. \hfill \textsc{Assertion}\\
%\textrm{⤳} 
\ex\label{ex:tom:26b}
Maria talked to at least a manager. \hfill \textsc{Presupposition}\\
%\textrm{⤳} 
\ex\label{ex:tom:26c}
Maria’s talking to a manager is not significant. \hfill \textsc{Presupposition}
\z 
\z 


\noindent Embedding under a question operator and under negation shows that the exclusive component in \REF{ex:tom:26a} is an assertion: it is targeted both by the question operator, \REF{ex:tom:27a}, and by the negation operator, \REF{ex:tom:28a}. %This is the exclusive component and we can show that this is an assertion by the fact that it is targeted by the question operator, \REF{ex:tom:27a}, and by the negation operator, \REF{ex:tom:28a}.
It is now readily apparent that the presupposition in \REF{ex:tom:26b} projects in a question, \REF{ex:tom:27b}, and under negation, \REF{ex:tom:28b}.
% Previous formulation:
% The presupposition in \REF{ex:tom:26b} is now readily apparent projecting in a question, \REF{ex:tom:27b}, and under negation, \REF{ex:tom:28b}.
The low-on-the-scale inference in \REF{ex:tom:26c} also clearly projects, \REF{ex:tom:27c}, \REF{ex:tom:28c}. 


\ea\label{ex:tom:27}
\gll Czy Maria rozmawiała tylko z menedżerem?\\  
Q Maria talked only with manager\\
\glt `Did Maria only/merely talk to a manager?’\\
%\textrm{⤳} 
\ea\label{ex:tom:27a}
Did Maria talk to anybody ranked higher than a manager?\hfill \textsc{Question}\\
%\textrm{⤳} 
\ex\label{ex:tom:27b}
Maria talked to at least a manager.\hfill \textsc{Presupposition}\\
%\textrm{⤳} 
\ex\label{ex:tom:27c} 
Maria’s talking to a manager is not significant.\hfill\textsc{Presupposition}
\z 
\z 

\ea\label{ex:tom:28}
\gll Maria nie rozmawiała tylko z menedżerem.\\  
Maria not talked only with manager\\
\glt `Maria did not merely talk to a manager.’\\
%\textrm{⤳} 
\ea\label{ex:tom:28a}
Maria talked to somebody ranked higher than a manager.  \\\hfill \textsc{Negative assertion}\\  
%\textrm{⤳} 
\ex\label{ex:tom:28b}
Maria talked to at least a manager. \hfill \textsc{Presupposition}\\
%\textrm{⤳} 
\ex\label{ex:tom:28c}
Maria’s talking to a manager is not significant. \hfill \textsc{Presupposition} 
\z 
\z

\noindent \citet{beaverclark08} and \citet{zeevat09} call the low-on-the-scale reading of \textit{only} \textsc{mirative}. As a mirative marker \textit{only} implies that the prejacent falls short of what is expected. This reading is easily detected with numerals, for instance, if you hear \textit{The driver is only 1 stop away}, you know that 1 stop is very close. However, the sentence \textit{The driver is only 15 stops away} also tells you that the driver is close -- contrary to what you would expect given the usual duration of 15 stops (so for this statement to be true these 15 stops must be very short). Without \textit{only}, the sentence \textit{The driver is 15 stops away} does not create any expectation whether the driver is going to reach you soon or not.


Consider also that \textit{tylko} can scope both below and above negation, as shown in \REF{ex:tom:29}. Both  readings can be detected: quantificational, \REF{ex:tom:29a}, \REF{ex:tom:29b}, and scalar, \REF{ex:tom:29c}, \REF{ex:tom:29d}.

\ea\label{ex:tom:29}
\gll Maria nie rozmawiała tylko z       menedżerem.\\
Maria not talked           only with manager\\
%\glt \textrm{⤳} 
\ea\label{ex:tom:29a}
‘Maria talked to somebody else in addition to the manager.’\\\hfill \textsc{not} > \textsc{only}\\
%\textrm{⤳} 
\ex\label{ex:tom:29b}
‘Maria talked to everyone else except the manager.’ \hfill \textsc{only} > \textsc{not}\\
%\textrm{⤳} 
\ex\label{ex:tom:29c}
‘Maria talked to somebody more noteworthy than the manager in addition to the manager.’ \hfill
\textsc{not} > \textsc{only}\\
%\textrm{⤳} 
\ex\label{ex:tom:29d}
‘Maria talked to everyone more noteworthy than the manager except the manager.’ \hfill \textsc{only} > \textsc{not}
\z 
\z 

\noindent Another test to distinguish between asserted and presupposed content are \textsc{emotive factive} predicates, such as `disappointed' in \REF{ex:tom:30} (\citealt{chierchiamcconnellginet90}, \citealt{beaverclark08}, \citealt{Simons.etal2011}, \citealt{coppockbeaver14}). These predicates target the asserted negative component with \textit{tylko}, showing, again, that exclusivity is asserted and scalarity presupposed, \REF{ex:tom:30}.

\ea\label{ex:tom:30}
\gll Jestem zawiedziony, że Maria rozmawiała tylko z menedżerem.\\
\textsc{aux} disappointed that Maria talked only with manager\\
\glt `I am disappointed that Maria did not talk to anybody else/more noteworthy.’
\z 

\noindent Example \REF{ex:tom:30} nicely illustrates the contribution of the prejacent presupposition ‘Maria talked to at least a manager’, because clearly the speaker is not disappointed that Maria did talk to the manager. A further test, utterance of assent/dis\-sent \citep{roberts11} is shown in \REF{ex:tom:31}: The exclusive component is asserted (targeted by direct affirmation) and the prejacent projects.

\ea\label{ex:tom:31}
\begin{xlist}
\exi{A:}{\gll Maria rozmawiała tylko z menedżerem.\\
Maria talked only with manager\\
\glt `Maria only talked to the manager.’}
\exi{B:}{\gll To prawda. \minsp{\{} Nie rozmawiała z nikim innym. / \minsp{\#} Rozmawiała z menedżerem.\}\\
this true {} not talked with nobody   else {} {} talked with manager\\
\glt `That is true. \{She did not talk to anybody else. / \textit{[Incongruous:]} She talked to the manager.\}’}
\exi{B$'$:}{\gll To prawda. \minsp{\{} Nie rozmawiała z kimś ważniejszym. /\hspace{.7cm} \minsp{\#} Rozmawiała z menedżerem.\}\\
this true {} not talked with somebody more.important {} {} talked with manager\\
\glt `That is true. \{She did not talk to somebody more noteworthy. / \textit{[Incongruous:]} She talked to the manager.\}’}
\end{xlist}
\z 


\noindent In the diagnostics above, \textit{tylko} patterns with \textit{only}, in that it has both scalar and quantificational readings and its exclusive component is asserted, while the prejacent and the low-on-the-scale components are presupposed. In fact, \textit{tylko} appears to be the most likely candidate for the counterpart of \textit{only}, since it is the only exclusive that can have quantificational readings (in contrast to \textit{zaledwie}, \textit{dopiero} and \textit{aż} discussed in the following sections). However, there is a significant difference between \textit{tylko} and \textit{only}: Although \textit{tylko} has scalar readings, it is incompatible with a temporal scale. The sentence in \REF{ex:tom:32a} can only mean that Marek woke up once (excluding all other ‘times of waking up’, i.e., the quantificational reading). The `no later than'-reading is available with \textit{(za)ledwie}, \textit{dopiero} and \textit{aż} but not with \textit{tylko} \REF{ex:tom:32b}.

\ea\label{ex:tom:32}
\ea\label{ex:tom:32a}      
\gll Martwi mnie, że Marek obudził się tylko o godzinie 13.\\
worries me that Marek woke self only at hour 13\\
\glt `I am worried that Marek woke up just once at 1 PM.’
\ex\label{ex:tom:32b}
\gll Martwi mnie, że Marek obudził się \minsp{\{} zaledwie / dopiero / aż /\hspace{.8cm} \minsp{\#} tylko\} o godzinie 13.\\
worries me that Marek woke self {} merely {} only\textsubscript{\textsc{temp}} {} \textsc{aż} {} {} only at hour 13\\
\glt `I am worried that Marek woke up only at 1 PM and no sooner.’
\z 
\z 

\noindent The contrast between  \REF{ex:tom:33a} and \REF{ex:tom:33b} shows that the presence of \textit{tylko} results in an anomaly, which indicates how strong the restriction on the type of scale is -- it is not possible to coerce a temporal reading.

\ea\label{ex:tom:33}
\ea[]{\label{ex:tom:33a} 
\gll Jest \minsp{\{} dopiero / zaledwie\} luty.\\
is {} only\textsubscript{\textsc{temp}} {} merely February\\
\glt `It is only February.’ (‘It is no later than February.’)}
\ex[\#]{\label{ex:tom:33b}             
\gll Jest tylko luty.\\
is only February\\
\glt \textit{Incongruous:} `It is no other month now than February.’}
\z 
\z 

\noindent If the dimension of the scale is underspecified in the semantics of an exclusive adverb, we expect it to be easily manipulated by the context (the semantics of the focused constituent the adverb associates with, e.g. `February' above, and the pragmatics). The four Polish exclusives, however, appear to be more contextually restricted. Distance, temporal and numeral scales are all natural, inherent, dimensional scales, so the prediction should be that an exclusive that works with any kind of a scale is available in any language. This seems to be the case for \textit{only} in English:

\ea\label{ex:tom:34}
We only reached Wilson Peak on that trip and were turned back [\dots] by an ice storm.~
\z 

\ea\label{ex:tom:35} I am usually up at 8:30am, but this morning for some reason I only woke up at 10am.
\z 

\ea\label{ex:tom:36} Brussels was ranked only 62\textsuperscript{nd} in Mercer's list of the world's most expensive cities for expats.
\z

\noindent If %the 
Polish \textit{tylko} can act scalar with individuals ordered on a contextual scale (importance, etc.) why is it incompatible with a temporal scale? In the next sections we will see that \textit{zaledwie} is dedicated to pragmatic scales (just like \textit{merely}) and \textit{dopiero} to temporal scales. This means that \textit{dopiero} might be blocking \textit{tylko} (but not \textit{zaledwie}) with a temporal scale, and that \textit{zaledwie} does not block \textit{tylko} with contextual scales (which perhaps suggests an interplay between the conventional nature of a scale and frequency in language use given that \textit{zaledwie} is less frequent than \textit{tylko}).

\subsection{\textit{(za)ledwie}}
\label{sec:tom:4.2}

While we saw in the previous section that \textit{tylko} has both quantificational and scalar readings, \textit{zaledwie} is specified for scalar readings -- it is a counterpart of %the 
English \textit{merely}. What makes it particularly interesting though, is that it is the exact scalar opposite of the exclusive \textit{aż}, which we will see in the next section. The example in \REF{ex:tom:37} shows that \textit{zaledwie} disprefers quantificational readings: In an impoverished context where no pragmatic ranking for Maria and others is available, the default is a scale of numbers of people who came (recall the two options for two scales in \REF{ex:tom:25} above) and so the use of \textit{tylko} is preferred.

\ea\label{ex:tom:37}
\gll \minsp{\{} Tylko / \minsp{?} Zaledwie\} Maria przyszła.\\
{} only {} {} merely  Maria came\\
\glt `Only/Merely Maria came.’\\
%\textrm{⤳} 
\ea \textsc{Quantificational reading}: `Nobody other than Maria came.'
\ex \textsc{Unavailable scalar reading}: `Nobody ranked higher than Maria came.'
\z 
\z 

\noindent Accordingly, \REF{ex:tom:38} shows that the most accessible interpretation for \textit{zaledwie} is that the excluded alternatives involve ranks higher than the manager (\ref{ex:tom:38}B$'$) and not people other than the manager (\ref{ex:tom:38}B).

\ea\label{ex:tom:38}
\begin{xlist}
\exi{A:}{\gll Maria rozmawiała zaledwie z menedżerem.\\
Maria talked merely with manager\\
\glt `Maria only talked to the manager.’}
%\ex\label{ex:tom:38b} 
\exi{B:}{\gll To prawda. \minsp{\{\#} Nie rozmawiała z innymi. / \minsp{\#} Rozmawiała z menedżerem.\}\\
this true {} not talked with others {} {} talked with manager\\
\glt `That is true. \{\textit{[Incongruous:]} She did not talk to other people. / \textit{[Incongruous:]} She talked to the manager.\}’}
%\ex\label{ex:tom:38c} 
\exi{B$'$:}{\gll To prawda. \minsp{\{} Nie rozmawiała z  kimś ważniejszym. /\hspace{0.8cm} \minsp{\#} Rozmawiała z menedżerem.\}\\
this true {} not talked with somebody more.important {} {} talked with manager\\
\glt `That is true. \{She did not talk to somebody more noteworthy. / \textit{[Incongruous:]} She talked to the manager.\}’}
\end{xlist}
\z 

\noindent Like \textit{tylko}, \textit{zaledwie} can scope both below and above negation, \REF{ex:tom:39}, which is expected given that \textit{tylko} can do so on both scalar and quantificational readings.

\ea\label{ex:tom:39}
\gll Maria nie rozmawiała zaledwie z menedżerem.\\
Maria not talked merely with manager\\
\ea 
%\textrm{⤳} 
‘Maria talked to somebody more noteworthy than the manager in addition to the manager.’\hfill \textsc{not} > \textsc{only}
%\textrm{⤳} 
\ex ‘Maria talked to everyone more noteworthy than the manager except the manager.’ \hfill \textsc{only} > \textsc{not}
\z 
\z 

\noindent The emotive factive diagnostic, just like with \textit{tylko} in \REF{ex:tom:28}, reveals the positive presupposition of the prejacent, \REF{ex:tom:40} (i.e., the speaker is not disappointed that Maria talked to the manager.)

\ea\label{ex:tom:40}
\gll Jestem zawiedziony, że Maria rozmawiała zaledwie z menedżerem.\\
 \textsc{aux} disappointed that Maria talked merely with manager\\
\glt `I am disappointed that Maria did not talk to anybody more noteworthy than the manager.’
\z 

\noindent What is interesting about \textit{zaledwie} is that of the four exclusive particles, it is the only one that can work as the counterpart of %the 
English \textit{barely} and receive \textsc{approximative} \REF{ex:tom:41} \citep{roberts11} and \textsc{counterfactual} \REF{ex:tom:42} \citep{toosarvandani10} readings. On approximative readings, \REF{ex:tom:41}, the prefix \textit{za}{}- can be dropped. On counterfactual readings, \REF{ex:tom:42}, \textit{ledwie} occurs without \textit{za-}.

\ea\label{ex:tom:41}
\gll \minsp{\{(} Za)ledwie / \minsp{\#} Tylko / \minsp{\#} Dopiero / \minsp{\#} Aż\} Marek otworzył  drzwi, a zadzwonił telefon.\\
{} merely {} {} only {} {} only\textsubscript{\textsc{temp}} {} {} \textsc{aż} Marek opened door and rang telephone\\
\glt `Marek barely opened the door, when the telephone rang.’
\z 

\ea\label{ex:tom:42}
\gll Marek \minsp{\{(*} za)ledwie / \minsp{\#} tylko / \minsp{\#} dopiero / \minsp{\#} aż\} otworzył drzwi.\\
Marek {} merely {} {} only {} {} only\textsubscript{\textsc{temp}} {} {} \textsc{aż} opened door \\     
\glt `Marek barely opened the door. (It was more likely for him to not open the door.)'
\z 

\noindent \textit{Ledwie} is syntactically an adverbial (also \textit{ledwo}), i.e. can only modify the VP, as shown by the contrast in \REF{ex:tom:43}: In \REF{ex:tom:43a} \textit{(za)ledwie} associates with Berlin, in \REF{ex:tom:43b} \textit{ledwie} associates with the VP ‘reached Berlin’ as can be seen from the translation, and here the prefix \textit{za}{}- is not possible.

\ea\label{ex:tom:43}
\ea\label{ex:tom:43a} 
\gll Dotarłem \minsp{(} za)ledwie do Berlina.\\
reached {} merely to  Berlin\\
\glt `I only reached Berlin.’
\ex\label{ex:tom:43b} 
\gll \minsp{(*} Za)ledwie dotarłem do Berlina.\\
 {} merely reached to Berlin\\
\glt `I barely made it to Berlin.’
\z 
\z 

\noindent Summing up, the exclusive \textit{zaledwie} prefers scalar readings (like \textit{merely} does), but its variant \textit{ledwie} has approximative and counterfactual readings.

\subsection{\textit{dopiero}}
\label{sec:tom:4.3}

In \REF{ex:tom:32a} above we saw that `It is only February' requires \textit{dopiero} in its translation, repeated below in \REF{ex:tom:44}, because \textit{tylko} results in an anomaly (`It is no other month now than February'). This suggests that \textit{dopiero} is a specialized \textsc{temporal} \textit{only} and I gloss it as \textit{only}\textsubscript{\textsc{temp}}.

\ea\label{ex:tom:44}
\ea\label{ex:tom:44a}   
It is only February.
\ex\label{ex:tom:44b}        
\gll Jest \minsp{\{} dopiero / \minsp{*} tylko\} luty.\\
is {} only\textsubscript{\textsc{temp}} {} {} only February\\
\glt `It is only February. (It is no later than February.)'
\z 
\z 

\noindent The use of \textit{dopiero}, in contrast to the other exclusives, coerces the scale associated with the predicate, e.g. quantity in \REF{ex:tom:45}--\REF{ex:tom:46}, to an interpretation with respect to time.

\ea\label{ex:tom:45}
\gll Maria rozmawiała dopiero z 2 osobami.\\
Maria talked only\textsubscript{\textsc{temp}} with 2 people\\
\glt `Maria only talked to 2 people so far.’
\z 

\noindent While in the example \REF{ex:tom:46} below temporal alternatives are obligatorily evoked, which results in the reading `so far Jan bought 2 tomatoes and he will buy more', there is no such expectation with \textit{tylko} in \REF{ex:tom:47}.

\ea\label{ex:tom:46}
\gll Jan kupił dopiero 2 pomidory.\\
Jan bought only\textsubscript{\textsc{temp}} 2 tomatoes\\
\glt `Jan bought only two tomatoes and we expect him to buy more.’
\z 

\ea\label{ex:tom:47}
\gll Jan kupił tylko 2 pomidory.\\
Jan bought only 2 tomatoes\\
\glt `Jan bought only two tomatoes.’
\z 

\noindent Logically, the temporal alternatives should be later in time: In \REF{ex:tom:45} more people will come later, in \REF{ex:tom:46} Jan will buy more tomatoes later, so in a context where such alternatives are not possible, as in \REF{ex:tom:48} where Jan is not going to write any more,  \textit{dopiero} is infelicitous.

\ea\label{ex:tom:48}
\gll Jan skończył pisać swój esej. Ma \minsp{\{} tylko / \minsp{\#} dopiero\} 10 stron.\\
Jan finished  writing his essay has {} only {} {} only\textsubscript{\textsc{temp}} 10 pages\\
\glt `Jan finished writing his essay. It is only 10 pages long.’
\z 

\noindent The example above provides an interesting contrast to the sentence in \REF{ex:tom:1} where the team ended up in 30th place (i.e., was not going to compete anymore) and \textit{dopiero} is felicitous. Indeed, we find \textit{dopiero} with scales that are not (necessarily) correlated with temporal progression, such as the example in \REF{ex:tom:printer-ex} about a partial achievement, idiomatically referred to as halfway to success.

\ea \label{ex:tom:printer-ex}   
\gll Dobra drukarka to jednak dopiero połowa sukcesu. Drugą, równie istotną częścią jest odpowiedni toner.\\ 
good printer is however only\textsubscript{\textsc{temp}} half         success second equally important part is suitable toner\\
\glt `A good printer is only half of the success. Another equally important thing is a good toner.’
\z 

\noindent Like the other exclusives we considered, \textit{dopiero} can take scope both above and below negation as we can see in \REF{ex:tom:49}--\REF{ex:tom:50}.\footnote{Native speakers of Polish can reflect on how additional context facilitates the interpretation of the two scope configurations. Consider also the different prosodic contours with different focus structures. In \REF{ex:tom:50a} \textit{dopiero} associates with focus on ‘2 people’, see \REF{ex:tom:fnia} below. In \REF{ex:tom:50b} \textit{dopiero} takes wide scope and associates with sentential focus, which can be made overt with scrambling as in \REF{ex:tom:fnib}.

\ea\label{ex:tom:fni} 
\ea\label{ex:tom:fnia} 
\gll Nie, nie przyszły dopiero \minsp{[} 2 osoby]\textsubscript{F}, są już \minsp{[} 3 osoby]\textsubscript{F}.\\
no not came only\textsubscript{\textsc{temp}} {} 2 people are already {} 3 people\\
\glt `So far not just 2 people showed up, but 3.’\hfill \textsc{not} > \textsc{only}
\ex\label{ex:tom:fnib} 
\gll Wydaje mi się, że \minsp{\{} nie przyszły dopiero 2 osoby / dopiero \minsp{[} 2 osoby nie przyszły]\textsubscript{F}\}, ale może się pomyliłam w liczeniu.\\
seems me self that {} not came only\textsubscript{\textsc{temp}} 2 people {} only\textsubscript{\textsc{temp}} {} 2 people not came but maybe self mistook in counting\\
\glt `I think that only 2 people didn't show up, but maybe I counted wrong’ \hfill \textsc{only} > \textsc{not}
\z 
\z
}

\ea\label{ex:tom:49}
\gll Dopiero 2 osoby przyszły.\\
Only\textsubscript{\textsc{temp}} 2 people came\\
\glt `So far only 2 people showed up.’
\z 

\ea\label{ex:tom:50}
\gll Nie przyszły dopiero 2 osoby.\\
not came only\textsubscript{\textsc{temp}} 2 people\\
%\glt \textrm{⤳} ?
\ea\label{ex:tom:50a} 
‘More than 2 people came.’   \hfill \textsc{not} > \textsc{only}\\
%\textrm{⤳} 
\ex\label{ex:tom:50b} 
‘Everyone came except 2 people.’ \hfill \textsc{only} > \textsc{not}
\z 
\z 

\noindent The translations for the examples above indicate that \textit{dopiero} is similar to the temporal adverbs \textit{so far} (\textit{jak} \textit{dotąd}) / \textit{yet} (\textit{jeszcze}). However, \textit{dopiero} also comes with an exclusive component which these temporal adverbs do not have by themselves. First, let us note that \textit{dopiero}, like \textit{yet}, and unlike \textit{so far}, can occur with negation and in questions, \REF{ex:tom:51}. Note also that  \textit{so far} is not an NPI in contrast to \textit{only} and \textit{yet}, \REF{ex:tom:51b}. 

\ea\label{ex:tom:51}
\gll Kto wie, czy to nie jest dopiero początek problemów premiera?\\
who knows if this not is only\textsubscript{\textsc{temp}} beginning problems prime.minister.\textsc{gen}\\
\glt `Who knows if this is only the beginning of the prime minister’s problems?’
\z 

\ea\label{ex:tom:51b}
Who knows, maybe this is \{only/yet/*so far\} the beginning of the prime minister’s problems?
\z

\noindent Secondly, with \textit{so far} in \REF{ex:tom:52a} there is an implication that Maria talked to just one person, but this is not an assertion because it is possible to continue the sentence as in \REF{ex:tom:52b} in contrast to \REF{ex:tom:52c} containing an exclusive. The counterpart of \REF{ex:tom:52c} in Polish, \REF{ex:tom:53}, is compatible with both \textit{tylko} and \textit{dopiero}.

\ea\label{ex:tom:52}
\ea\label{ex:tom:52a}  
So far Maria has talked to the manager.
\ex\label{ex:tom:52b}        
So far Maria has talked to the manager and a couple of other people.
\ex\label{ex:tom:52c}       
So far Maria has only talked to the manager \#and a couple of other people.
\z 
\z 

\ea\label{ex:tom:53}
\gll Jak dotąd Maria rozmawiała \minsp{\{} dopiero / tylko\} z menedżerem.\\
how till.now Maria talked {} only\textsubscript{\textsc{temp}} {} only with manager\\
\glt `So far Maria has only talked to the manager.’ 
\z 

\noindent That \textit{so far} does not assert the exclusion of alternatives is also apparent in \REF{ex:tom:54}--\REF{ex:tom:56} below. In \REF{ex:tom:54} we can expect that the person will call later, and accordingly in the Polish translation in \REF{ex:tom:55} \textit{dopiero} is not felicitous.

\ea\label{ex:tom:54}
So far he has not called.
\z 

\ea\label{ex:tom:55}
\gll \minsp{\{} Jak dotąd / \minsp{*} dopiero\} nie zadzwonił.\\
{} how till.now {} {} only\textsubscript{\textsc{temp}} not called\\
\glt ‘So far he has not called.' (Lit. ‘Up till now he has not called.’)
\z

\noindent The presence of \textit{so far} and \textit{only} in \REF{ex:tom:56} together results in an implication that documenting 373 species is a lot, because ‘it is still only February’, as said in 2012, indicates that this number was reached earlier than expected.\footnote{The source of \REF{ex:tom:56} is \url{https://james-champion.com/friday-24th-february-2012/} (accessed on 29 March 2012). In fact, John’s initial goal for 2012 was 500 species and by the end of the year he documented 591 species.}

\ea\label{ex:tom:56}
John Cahill is going for an all-out mega-bird list for Guatemala in 2012. [\dots]. So far he has reached 373 species in Guatemala this year -- and it is still only February!
\z

\noindent In contrast to \textit{so far} in \REF{ex:tom:56}, \textit{dopiero} in the sentence in \REF{ex:tom:57} contributes the reading that 373 species is not much, because, as the utterance of assent test shows, \REF{ex:tom:58}, the exclusive component ‘no more than 373 yet’  is asserted. The same test shows that \textit{jak dotąd} (\textit{so far}) lacks the exclusive component.

\ea\label{ex:tom:57}
\begin{xlist}
\exi{A:}{\gll Jan dotarł dopiero do 373 gatunków.\\
Jan reached only\textsubscript{\textsc{temp}} to 373 species\\
\glt ‘Jan has only reached 373 species so far.'}
\exi{B:}{\gll To prawda. \minsp{\{} Nie dotarł do więcej niż 373 gatunków. / \minsp{\#} Dotarł do 373 gatunków.\}\\
this true {} not reached to more than 373 species {} {} reached to 373 species\\
\glt ‘That's true. \{He hasn't reached more than 373 species. / \textit{[Incongruous:]} He has reached 373 species.\}'}
\end{xlist}
\z

\ea\label{ex:tom:58}
\begin{xlist}
\exi{A:}{\gll Jan dotarł jak dotąd do 373 gatunków.\\
Jan reached so far to 373 species.\\
\glt ‘Jan has so far reached 373 species.'}
\exi{B:}{\gll To prawda. Dotarł \minsp{(} już) do 373 gatunków.\\
this true reached {} already to 373 species\\
\glt ‘That's true. He has (already) reached 373 species.'}
\end{xlist}
\z

\largerpage
\noindent The data above demonstrates that \textit{dopiero} specializes for temporal scales and now we will see that the diagnostics for assertion and presupposition replicate the results we found for the other exclusives. The utterance of assent test, \REF{ex:tom:59}, the emotive factive predicate test, \REF{ex:tom:60} and negation test, \REF{ex:tom:61}, show that the exclusion of alternatives higher on the (temporal) scale is asserted, the prejacent is presupposed and the low-on-scale reading is presupposed.
%excludes alternatives ordered along a temporal scale, therefore I call it the \textsc{temporal} \textit{only} (only\textsubscript{\textsc{temp}}), which can be translated into English as ‘\textit{so far only}’. Using \textit{dopiero} in the same contexts as I used earlier to diagnose assertion and presupposition (utterance of dissent test \REF{ex:tom:60}, the emotive factive predicate test in \REF{ex:tom:61} and negation in \REF{ex:tom:62}) we replicate the results for the other exclusives: The exclusion of alternatives higher on the (temporal) scale is asserted, the prejacent is presupposed and the low-on-scale reading is presupposed.

\ea\label{ex:tom:59}
\begin{xlist}
\exi{A:}{\gll Maria rozmawiała dopiero z menedżerem.\\
Maria talked only\textsubscript{\textsc{temp}} with manager\\
\glt `Maria only talked to the manager so far.’}
\exi{B:}{\gll To prawda. \minsp{\{} Nie rozmawiała \minsp{(} jeszcze) z innymi. / \minsp{\#} Rozmawiała z menedżerem.\}\\
this true {} not talked {} yet with others {} {} talked with manager\\
\glt `That is true. \{She has not talked to other people yet. / \textit{[Incongruous:]} She talked to the manager.\}’}
\exi{B$'$:}{\gll To prawda. \minsp{\{} Nie rozmawiała \minsp{(} jeszcze) z nikim {ważniejszym. /} \minsp{\#} Rozmawiała z menedżerem.\}\\
this true {} not she.talked {} yet with nobody more.important {} she.talked with manager\\
\glt `That is true. \{She has not talked to anybody more noteworthy yet. / \textit{[Incongruous:]} She talked to the manager.\}’}
\end{xlist}
\z 

\ea\label{ex:tom:60}
\gll Jestem zawiedziony, że Maria rozmawiała dopiero z menedżerem.\\
 \textsc{aux} disappointed that Maria talked only\textsubscript{\textsc{temp}}   with manager\\
\glt `I am disappointed that Maria hasn’t talked to anybody else/more noteworthy yet.’ (\textit{Implied:} ‘Maria talked to the manager.’)\\
%$\rightsquigarrow$ ‘Maria talked to the manager.’
\z 

\ea\label{ex:tom:61}
\gll Maria nie rozmawiała dopiero z menedżerem.\\
Maria not talked only\textsubscript{\textsc{temp}} with manager\\
%\glt \textrm{⤳} 
\ea
‘Maria has already talked to somebody more noteworthy than the manager in addition to the manager.’ \hfill \textsc{not} > \textsc{only}\\
%\textrm{⤳} 
\ex
‘Maria has already talked to everyone more noteworthy than the manager except the manager.’ \hfill  \textsc{only} > \textsc{not}
\z 
\z 

\noindent Finally, we must note that while \textit{dopiero} excludes stronger alternatives on a temporal scale, they are not necessarily later, as is clear in \REF{ex:tom:62}. What is excluded in \REF{ex:tom:62a} are earlier alternatives for finishing the writing. In this context, the exclusive that I am going to present in the next section, \textit{aż}, is also possible. The fact that in \REF{ex:tom:62b} \textit{aż} is not felicitous, but instead it requires the preposition \textit{do}, lit. ‘to’, \REF{ex:tom:62c}, which carries the meaning of `until' (i.e., here \textit{aż} is optional but \textit{do} is not; see also \REF{ex:tom:plfriday} in \sectref{sec:tom:6}) indicates that \textit{aż} is not inherently specified for temporal scales and what allows it in \REF{ex:tom:62a} is the presence of the numeral.

\ea\label{ex:tom:62}
\ea\label{ex:tom:62a} 
\gll Skończę pisać \minsp{\{} dopiero / aż / \minsp{\#} tylko / \minsp{\#} zaledwie\} za 2 miesiące.\\
will.finish writing {} only\textsubscript{\textsc{temp}} {} \textsc{aż} {} {} only {} {} merely in 2 months\\
\glt ‘I will finish writing no sooner than in 2 months.’
\ex\label{ex:tom:62b}            
\gll Skończę pisać \minsp{\{} dopiero / \minsp{\#} aż / \minsp{\#} tylko / \minsp{\#} zaledwie\} w sobotę.\\
will.finish writing {} only\textsubscript{\textsc{temp}} {} {} \textsc{aż} {} {} only {} {} merely on Saturday\\
\glt ‘I will finish writing no sooner than on Saturday.’
\ex\label{ex:tom:62c}            
\gll Skończę pisać \minsp{\{} aż / \minsp{\#} dopiero / \minsp{\#} tylko / \minsp{\#} zaledwie\} do soboty.\\
will.finish writing {} \textsc{aż} {} {} only\textsubscript{\textsc{temp}} {} {} only {} {} merely until Saturday\\
\glt ‘I will finish writing until Saturday, which is late.’
\z 
\z 

\noindent We can conclude that \textit{dopiero} has the same meaning components as \textit{tylko}/\textit{zale\-dwie}, but it is specified for numeral, especially temporal scales. In the next section, I will show that \textit{aż} is a scalar opposite of \textit{dopiero}.

\subsection{\textit{aż}}
\label{sec:tom:4.5}

The exclusive \textit{aż} is a scalar opposite of \textit{zaledwie} (\citealt{tomaszewicz12,tomaszewicz13,tomaszewicz13a}). The assent test in \REF{ex:tom:63} shows us first that \textit{aż} has no quantificational reading. Neither of the continuations in (\ref{ex:tom:63}B) are possible -- just like with \textit{zaledwie} in (\ref{ex:tom:38}B). Secondly, we see in (\ref{ex:tom:63}B$'$) that the asserted component is `She did not talk to somebody less noteworthy', which is the direct opposite of what we saw in (\ref{ex:tom:38}B$'$) with \textit{zaledwie}: `She did not talk to somebody more noteworthy'. 
%(In English, there is no exclusive that would be the scalar opposite of \textit{merely}, hence there is no gloss for \textit{aż}.)


\ea\label{ex:tom:63}
\begin{xlist}
\exi{A:}{\gll Maria rozmawiała aż z menedżerem.\\
Maria talked \textsc{aż} with manager\\
\glt `Maria talked to somebody so important as the manager.’}
\exi{B:}{\gll To prawda. \minsp{\{\#} Nie rozmawiała z nikim innym. / \minsp{\#} Rozmawiała z menedżerem\}.\\
this true {} not talked with nobody other {} {} talked with manager\\
\glt `That is true. \{\textit{[Incongruous:]} She did not talk to other people. / \textit{[Incongruous:]} She talked to the manager.\}'}
\exi{B$'$:}{\gll To prawda. \minsp{\{} Nie rozmawiała z kimś mniej ważnym. /\hspace{.7cm} \minsp{\#} Rozmawiała z menedżerem.\}\\
this true {} not talked with somebody less important {} {} talked with manager\\
\glt `That is true. \{She did not talk to somebody less noteworthy. / \textit{[Incongruous:]} She talked to the manager.\}’}
\end{xlist}
\z 

\noindent Native speaker intuitions are particularly robust for the juxtaposition of \textit{aż} and \textit{zaledwie} in \REF{ex:tom:64}--\REF{ex:tom:67} where the reason clause identifies the asserted content (\citealt{dretske72}, \citealt{beaverclark08}). In \REF{ex:tom:64}--\REF{ex:tom:65} the reason for Maria’s (un)hap\-pi\-ness is the position of the manager on the scale: high with \textit{aż}, low with \textit{zaledwie}. In \REF{ex:tom:66}--\REF{ex:tom:67} the reason for the party’s success/failure is the fact that the expected number of guests was either exceeded, \REF{ex:tom:66}, or not achieved, \REF{ex:tom:67}. With \textit{aż} in \REF{ex:tom:64} Maria is happy because she is not talking to somebody less important than the manager and in \REF{ex:tom:66} the party was a success because no less than 15 people came. The mere fact that Maria is talking to the manager is not the reason for her happiness -- the fact that the manager is high on the scale of the relevant people is the reason. Similarly, the fact that 15 people came to the party is not the reason for the success -- the fact that this number exceeded the expectations is the reason. Conversely, with \textit{zaledwie} in \REF{ex:tom:65} Maria is unhappy because her goal is to talk to someone more important than the manager and in \REF{ex:tom:67} the party was a failure because more than 15 people were expected. The reason clause test demonstrates that exclusivity is asserted, which means that \textit{aż} works like the scalar opposite of \textit{zaledwie}, \textit{tylko}  and \textit{dopiero}: It carries a high-on-the-scale presupposition.

\ea\label{ex:tom:64}
\gll Maria jest zadowolona, bo rozmawia, aż z menedżerem.\\
Maria is happy because talks \textsc{aż} with manager\\
\glt `Maria is happy because she is not talking to somebody less important than the manager.’\\
%\textit{Unavailable:} ‘Maria is happy because she is talking to the manager.’
\z 

\ea\label{ex:tom:65}
\gll Maria jest niezadowolona, bo rozmawia zaledwie z menedżerem.\\
Maria is unhappy because talks merely with manager\\
\glt `Maria is unhappy because she is not talking to somebody more important than the manager.’\\
%\textit{Unavailable:} ‘Maria is unhappy because she is talking to the manager.’
\z 

\ea\label{ex:tom:66}
\gll Impreza się udała, bo przyszło aż 15 osób.\\
party self succeeded because came \textsc{aż} 15 people\\
\glt `The party was a success because no less than 15 people came.’\\
%\textit{Unavailable:} ‘The party was a success because 15 people came.’
\z 

\ea\label{ex:tom:67}
\gll Impreza się nie udała, bo  przyszło \minsp{\{} tylko / zaledwie\} 15 osób.\\
party self not succeeded because came {} only {} merely 15 people\\
\glt `The party was not a success because no more than 15 people came.’\\
%\textit{Unavailable:} ‘The party was not a success because 15 people came.’
\z 

\noindent A parallel example with \textit{dopiero} does not work, because it evokes a temporal alternative when more people come.

\ea\label{ex:tom:party-dopiero}   
\gll Impreza się \minsp{(} nie) udała, \minsp{\#} bo przyszło dopiero 15 osób.\\
party self {} not succeeded {} because came only\textsubscript{\textsc{temp}} 15 people\\
\glt `The party was (not) a success \textit{[Incongruous:]} because so far only 15 people came.’
\z 

\noindent The high-on-the-scale presupposition found with \textit{aż} projects under negation and in questions, and the orientation of the scale stays the same, as shown in \REF{ex:tom:68c}, \REF{ex:tom:69c}, \REF{ex:tom:70c}. Embedding under question and negation operators clearly reveals the asserted exclusive component -- the scalar opposite of the contribution of \textit{only}. Although \REF{ex:tom:68} could simply be taken to convey that Maria talked to a manager and that this is significant, its negative version, \REF{ex:tom:69}, crucially asserts not only that Maria did not talk to a manager, but also that she talked to somebody lower on the scale than the manager, \REF{ex:tom:69a}. This is evidence for an asserted exclusive component.

\ea\label{ex:tom:68}
\gll Maria rozmawiała aż z  menedżerem. \\         
Maria talked \textsc{aż} with  manager\\
\glt `Maria talked to somebody so important as the manager.’\\
%\textrm{⤳} (a) 
\ea\label{ex:tom:68a} 
Maria talked to nobody ranked lower than a manager.\hfill \textsc{Assertion}\\
%\textrm{⤳} (b) 
\ex\label{ex:tom:68b} 
Maria talked to at most a manager.  \hfill    \textsc{Presupposition}\\
%\textrm{⤳} (c) 
\ex\label{ex:tom:68c} 
Maria’s talking to a manager is significant. \hfill \textsc{Presupposition}
\z 
\z 

\ea\label{ex:tom:69}
\gll Maria nie rozmawiała aż z menedżerem. \\     
Maria not talked \textsc{aż} with manager\\
\glt `Maria talked to somebody less important than the manager.’\\
%\textrm{⤳} (a) 
\ea\label{ex:tom:69a}
Maria talked to somebody ranked lower than a manager. \hfill \textsc{Assertion}\\
%\textrm{⤳} (b) 
\ex\label{ex:tom:69b} 
Maria talked to at most a manager. \hfill \textsc{Presupposition}\\
%\textrm{⤳} (c) 
\ex\label{ex:tom:69c} 
Maria’s talking to a manager is significant. \hfill \textsc{Presupposition}
\z 
\z 

\noindent Exclusivity also emerges in the question in \REF{ex:tom:70} which asks whether the position of the person Maria talked to is lower than a manager, \REF{ex:tom:70a}. (This is exactly the opposite of the effect of the presence of scalar \textit{tylko} `only' in the question in \REF{ex:tom:28}).

\ea\label{ex:tom:70}
\gll Czy Maria rozmawiała aż z menedżerem?  \\  
Q Maria talked \textsc{aż} with manager\\
\glt `Did Maria talk to somebody so important as the manager?’\\
%\textrm{⤳} (a) 
\ea\label{ex:tom:70a} 
Did Maria talk to somebody ranked higher than a manager?\\\hfill \textsc{Question}\\
%\textrm{⤳} (b) 
\ex\label{ex:tom:70b} 
Maria talked to at most a manager.\hfill \textsc{Presupposition}\\
%\textrm{⤳} (c) 
\ex\label{ex:tom:70c} 
Maria’s talking to a manager is significant. \hfill \textsc{Presupposition}
\z 
\z 

%\noindent The exclusive component of \textit{aż} in the a.-examples of \REF{ex:tom:68}--\REF{ex:tom:70}, i.e. the assertion that the alternatives lower on the scale are excluded, is the reverse of the contribution of \textit{only}/\textit{merely}, where higher alternatives are excluded. The two other meaning components of \textit{aż} that project in the b./c.-examples \REF{ex:tom:68}--\REF{ex:tom:70} are scalar reversals of the two presuppositions of \textit{only/merely} in the b./c.-examples of \REF{ex:tom:26}--\REF{ex:tom:28}.

\noindent The exclusive component in \REF{ex:tom:68a}, \REF{ex:tom:69a} and \REF{ex:tom:70a}, i.e. the assertion that the alternatives lower on the scale are excluded, is the reverse of the contribution of \textit{only}/\textit{merely}, where higher alternatives are excluded, as we saw in \REF{ex:tom:26a}, \REF{ex:tom:27a} and \REF{ex:tom:28a}. The two other meaning components of \textit{aż} that project are scalar reversals of the two presuppositions of \textit{only/merely} identified in (\ref{ex:tom:26}b/c), (\ref{ex:tom:27}b/c) and (\ref{ex:tom:28}b/c).


%In \REF{ex:tom:68}--\REF{ex:tom:70}, the presupposition in b. is that the alternatives under consideration can be at most as high on the scale as the prejacent. The presupposition in c. is that the prejacent is significantly high on the scale. For \REF{ex:tom:68} this gives us the interpretation that considering the alternative positions that the person Maria talked to could hold \REF{ex:tom:68b}, being a manager is the most significant among them \REF{ex:tom:68b}--\REF{ex:tom:68c}, and that that person is indeed no less than a manager \REF{ex:tom:68a}.

In \REF{ex:tom:68b}, \REF{ex:tom:69b}, \REF{ex:tom:70b}, the presupposition is that the alternatives under consideration can be at most as high on the scale as the prejacent. The presupposition in \REF{ex:tom:68c}, \REF{ex:tom:69c}, \REF{ex:tom:70c} is that the prejacent is significantly high on the scale. For \REF{ex:tom:68} this gives us the interpretation that considering the alternative positions that the person Maria talked to could hold \REF{ex:tom:68b}, being a manager is the most significant among them \REF{ex:tom:68b}, \REF{ex:tom:68c}, and that that person is indeed no less than a manager \REF{ex:tom:68a}.

\citet{beaverclark08} sum up the contribution of \textit{only} as “contra expectation, nothing stronger holds” (p. 279). The scalar reversal of each of the meaning components of \textit{only}, yields the interpretation for \textit{aż} that can be described as `contra expectation, nothing weaker holds'.

It is frequently noticed that \textit{even} functions as an antonym of \textit{only} at the level of presupposition (\citealt{horn69}, \citealt{koenig91}, \citealt{beaverclark08}, \citealt{greenberg22}). The following contrasts between \textit{aż}, \textit{zaledwie} and \textit{nawet} `even' show that the scalar presupposition of \textit{aż} and \textit{zaledwie} picks up the contextually salient scale, namely, distance, while \textit{even} requires the likelihood scale (contra the speculation about \textit{aż} vs. \textit{nawet} in \citealt{bajaj16}).~Negation in the presence of \textit{aż} in \REF{ex:tom:71} conveys that the distance that Jan traveled was shorter than to Warsaw; the sentence cannot mean that Jan did not travel at all, which is a possible interpretation in \REF{ex:tom:12} containing negation and \textit{even}. The sentence in \REF{ex:tom:71} tells us both that (i) Jan went as far as Łódź, but not Warsaw (see Figure \ref{fig:tom:warsaw} below for the illustration of the distance), and that (ii) Jan’s trip was not noteworthy, because it was not far. This meaning results from negation scoping above~\textit{aż} and is exactly parallel to the meaning of \REF{ex:tom:73} where negation scopes above \textit{zaledwie}. \REF{ex:tom:73} says that (i) Jan went further than Łódź, Jan went as far as Warsaw, and that (ii) Jan’s trip was noteworthy, because it was far. Both the use of \textit{aż} and \textit{zaledwie} induce a distance reading in contrast to \textit{nawet} `even' in \REF{ex:tom:72}. In \REF{ex:tom:72} \textit{nawet} needs to scope above negation, the scalar presupposition is that the prejacent `Jan went to Łódź' is the most likely among the alternatives (i.e., the least significant).

\ea\label{ex:tom:71}
\gll Jan nie pojechał aż do Warszawy, \minsp{\{} pojechał zaledwie do Łodzi / \minsp{\#} został w domu\}.\\             
Jan not drove \textsc{aż} to Warsaw {} drove merely to Łódź {} {} stayed  at  home\\
\glt `Jan did not drive as far as Warsaw, \{he drove to Łódź / \textit{[Incongruous:]} he stayed at home\}.’
\z 

\ea\label{ex:tom:72}
\gll Jan nie pojechał nawet do Łodzi, \minsp{\{\#} pojechał zaledwie do Warszawy. / został w domu\}. \\                   
Jan not drove even to Łódź {} drove merely to Warsaw {} stayed  at  home\\
\glt `Jan did not even drive to Łódź, \{\textit{[Incongruous:]} he drove merely to Warsaw / he stayed at home\}.’
\z

\ea\label{ex:tom:73}
\gll Jan nie pojechał zaledwie do Łodzi, \minsp{\{} pojechał aż do Warszawy. / \minsp{\#} został w   domu\}. \\     
Jan not drove merely to Łódź {} drove \textsc{aż} to Warsaw {} {} stayed at  home\\
\glt `Jan did not go merely to Łódź, \{he went all the way to Warsaw / \textit{[Incongruous:]} he stayed at home\}.’
\z

\begin{figure}
    \centering
    % 
\includegraphics[height=0.2\textheight]{figures/az_map_crop.png}%
    \includegraphics[height=0.2\textheight]{figures/osm_map.png}
    \caption{An illustration of the distances in the situation in example \REF{ex:tom:71}: ‘Jan did not go as far as Warsaw, he went merely to Łódź.’\\{\tiny © OpenStreetMap contributors, Open Database License}}
    \label{fig:tom:warsaw}
\end{figure}

%%please move the includegraphics inside the {figure} environment
%%\includegraphics[width=\textwidth]{figures/AZSlavicSemanticsFINAL-img001.png} & \\

In addition to its use as a scalar particle, \textit{aż} functions in Polish as a complementizer, equivalent to \textit{until} in English, \REF{ex:tom:74a}. As a complementizer, it is obligatorily sentence-initial, but as we saw in \sectref{sec:tom:3} focus associators can associate with foci of all sizes, including the whole sentence in Polish. I propose that the lexical semantics of \textit{aż} is the same when it is interpreted as \textit{until}, \REF{ex:tom:74a}--\REF{ex:tom:74b}, and when it is a scalar adverb, \REF{ex:tom:74c}, but the syntax determines whether the ordering of alternatives is on a temporal scale or a pragmatic scale.\footnote{\label{fn:tom:former4}I thank Radek Šimík for pointing out that the other exclusive particles cannot switch in a parallel way with respect to the syntax, suggesting that the semantics of \textit{aż} alone cannot be responsible for its function as a complementizer. This is an interesting topic for (a cross-linguistic) investigation.} I observe that when the main clause explicitly mentions duration, \textit{aż} can either work as \textit{until}, \REF{ex:tom:74b}, or as the scalar exclusive in which case it is preceded by the complementizer \textit{that}, \REF{ex:tom:74c}. In \REF{ex:tom:74c} \textit{aż} can either associate with the subject so that the alternatives are people who could tell Ann to stop singing, or with the whole clause so that alternatives are other less serious things that may happen, e.g., the whole auditioning committee laughing.

\ea\label{ex:tom:74}
\ea\label{ex:tom:74a} 
\gll Anna śpiewała, aż  [\textsubscript{IP}[\textsubscript{DP} Maria] kazała jej przestać]. \\ 
Anna sang \textsc{aż} {} Maria ordered her stop\\
\glt `Anna sang until Maria told her to stop.’
\ex\label{ex:tom:74b}       
\gll Anna śpiewała tak długo, aż  [\textsubscript{IP}[\textsubscript{DP} Maria] kazała jej przestać]. \\ 
Anna sang so long \textsc{aż} {} Maria ordered her stop\\
\glt `Anna kept on singing until Maria told her to stop.’
\ex\label{ex:tom:74c}       
\gll Anna śpiewała tak \minsp{\{} słabo / długo\}, że aż [\textsubscript{IP}[\textsubscript{DP} Maria] kazała jej przestać]. \\ 
Anna sang so {} badly {} long that \textsc{aż} {} Maria ordered her stop\\
\ea `Anna sang so badly/long, that out of all things that could happen Maria told her to stop.’\\
\ex `Anna sang so badly/long, that out of all people Maria told her to stop.’
\z 
\z 
\z 

\noindent In this section we established that \textit{aż} is a scalar opposite of \textit{zaledwie}, scalar \textit{tylko} and \textit{dopiero.} Similarly to \textit{dopiero}, it receives a temporal interpretation when in the complementizer position. In the next section, I show how the lexical entries for the four exclusives can be related by a common core semantics.

\section{Analysis}
\label{sec:tom:5}

I propose that the four exclusives, \textit{tylko}, \textit{zaledwie}, \textit{dopiero} and \textit{aż} are a family related by having the same three semantic components: the exclusive assertion, the scalar presupposition and the prejacent/factive presupposition. The differences between them involve the direction of the scale and extra requirements on the sets of alternatives resulting in their idiosyncrasies. The facts are summarized in \tabref{tab:tom:1}. Below I define the lexical entries for each member of the family in such a way that the common core of the three components is present and only the necessary parameters are tweaked.

% \begin{table}[]
%     \centering
%     \begin{tabular}{lcccc}
%     \lsptoprule
%     &\textit{tylko}&\textit{(za)ledwie} & \textit{dopiero} & \textit{aż}  \\
%     \midrule
%     prejacent $p$:&\multicolumn{4}{c}{$p$ is presupposed}\\ 
%     scalar pre-&\multicolumn{3}{c}{$p$ is low on the scale $S$}&$p$ is high\\ 
%     supposition:&\multicolumn{3}{c}{}&on the scale $S$\\ 
%     exclusive&\multicolumn{3}{c}{alternatives higher than $p$ on $S$ are excluded}&alternatives lower\\ 
%     assertion:&\multicolumn{3}{c}{}&than $p$ on $S$\\ 
%     &\multicolumn{3}{c}{}&are excluded\\ 
%     idiosyncrasies:&$S$ is logical or&$S$ is pragmatic;&$S$ is numeral/&$S$ is pragmatic;\\ 
%     &pragmatic but&approximative&temporal&temporal `until'\\ 
%     &not temporal&readings&&meaning\\ 
%     \lspbottomrule
%     \end{tabular}
%     \caption{The meaning components of \textit{tylko}, \textit{zaledwie}, \textit{dopiero} and \textit{aż}}
%     \label{tab:tom:sum}
% \end{table}

\begin{table}
\begin{tabularx}{\textwidth}{p{0.2cm}XXXX}
\lsptoprule
 & \textit{tylko} & \textit{(za)ledwie} & \textit{dopiero} & \textit{aż}\\\midrule
\multicolumn{4}{l}{\textsc{status of prejacent} $p$}&{}\\
&\cellcolor{black!10}\textit{p} presupposed&\cellcolor{black!10}\textit{p} presupposed&\cellcolor{black!10}\textit{p} presupposed&\cellcolor{black!10}\textit{p} presupposed\\[1ex]
\multicolumn{5}{l}{\textsc{scalar presupposition ($p$ high/low on some scale} $S$)}\\
&\cellcolor{black!10}\textit{p} low on \textit{S} &\cellcolor{black!10}\textit{p} low on \textit{S}&\cellcolor{black!10}\textit{p} low on \textit{S}&\cellcolor{black!20}\textit{p} high on \textit{S}\\[1ex]
\multicolumn{5}{l}{\textsc{exclusive assertion (alternatives excluded, higher/lower on $S$ than $p$)}}\\
&\cellcolor{black!10}higher than \textit{p}&\cellcolor{black!10}higher than \textit{p}&\cellcolor{black!10}higher than \textit{p}&\cellcolor{black!20}lower than \textit{p}\\[1ex]
\textsc{idiosyncrasies} &\\
&\cellcolor{black!10}\textit{S} is logical or pragmatic but not temporal&\cellcolor{black!20}\textit{S} is pragmatic; approximative readings &\cellcolor{black!10}\textit{S} is numeral/ temporal&\cellcolor{black!20}\textit{S} is pragmatic; temporal `until' meaning\\
\lspbottomrule
\end{tabularx}
\caption{The meaning components of \textit{tylko}, \textit{zaledwie}, \textit{dopiero} and \textit{aż}}
\label{tab:tom:1}
\end{table}

The meaning of the scalar exclusive, \textit{tylko/zaledwie} involves the assertion that no higher alternative is true, and the factive presupposition (\ref{ex:tom:75a}), as well as the low-on-the-scale presupposition (\ref{ex:tom:75b}). Scalarity is thus both asserted and presupposed.

% \ea\label{ex:tom:75}
% \sib{\textit{tylko/zaledwie}}$^w (C)(p)$ is defined iff\\
% (i) $\exists q \in C [ q \ge p \wedge q(w) = 1]$ \\\hfill \textit{Factivity: an alternative at least as strong as $p$ is true}\\
% and\\
% (ii) $\forall q \in C [ q \ne p \rightarrow q > p]$ \hfill \textit{Scalarity: $p$ is the weakest alternative}\\ 
% if defined then \sib{\textit{tylko/zaledwie}}$^w (C)(p) = 
% \forall q [ q \in C \wedge q \ne p \wedge q > p)\rightarrow q(w) = 0]$
% \z 

\ea\label{ex:tom:75}
\sib{\textit{tylko/zaledwie}}$^w(C)(p)$ is defined iff
\ea $\exists q\in C [q\geq p\wedge q(w) = 1]$\label{ex:tom:75a}\\
{[Factivity: an alternative at least as strong as $p$ is true]}\smallskip\\
and
\ex ${\forall} q\in C [q\neq p → q > p]$\label{ex:tom:75b}\\
{[Scalarity: $p$ is the weakest alternative]}\smallskip
\z
if defined then\\
\sib{\textit{tylko/zaledwie}}$^w(C)(p)= \forall q[q\in C\wedge q\neq p\wedge q > p) → q (w) = 0]$ 
\z

The \textsc{factive} presupposition is the positive meaning component that projects under negation, therefore it is not just the prejacent proposition, but an alternative proposition that is at least as strong on the relevant scale (I am following \citealt{klinedinst05} and \citealt{coppockbeaver14}, but see also \citealt{beaverclark08} and \citealt{zeevat09} for a mirative presupposition that the prejacent falls short of what is expected). This positive contribution is paraphrasable with `at least' as we first saw in the example in \REF{ex:tom:16}: `Maria only talked to the manager' implies that `Maria talked to (at least) the manager', (\ref{ex:tom:16}c). \citet{coppockbeaver14} point out that this component is easier to see on scalar readings, in examples such as \REF{ex:tom:coppockbeaver14a}, than on quantificational readings, as in \REF{ex:tom:coppockbeaver14b}.
%e.g., `This is not \{just/a mere\} down payment' clearly tells us that `This is at least a down payment' (their examples (8--9) on p. 374). Their example for the quantificational reading is the sentence `Mary did not invite only John and Mike' which implies that `Mary invited at least John and Mike' ((18) on p. 379). 

\ea\label{ex:tom:coppockbeaver14a}
This is not \{just/a mere\} down payment.\smallskip\\
\textit{Positive contribution}:
`This is at least a down payment.'\\\hfill \textsc{scalar reading} \cite[(8--9) p. 374]{coppockbeaver14} 
\z

\ea\label{ex:tom:coppockbeaver14b}  
Mary did not invite only John and Mike. \smallskip\\
\textit{Positive contribution}:
`Mary invited at least John and Mike.'\\\hfill \textsc{quantificational reading} \cite[(18) p. 379]{coppockbeaver14}
\z


Since my analysis follows the accounts where the quantificational and scalar readings are derived from the same lexical entry (\citealt{bonomicasalegno93}, \citealt{vanrooij02}, \citealt{beaver04}, \citealt{klinedinst05}, \citealt{riester06}, \citealt{beaverclark08}, \citealt{coppockbeaver10, coppockbeaver14}), the factive presupposition of \textit{tylko/zaledwie} in (\ref{ex:tom:75a}) allows for both readings. As mentioned in \sectref{sec:tom:2}, the dimension of the scale is set by the context and  quantificational readings arise with logical scales, whereas scalar readings result from pragmatic orderings of alternatives in \textit{C.} In impoverished contexts the logical scale is the default.

\citet{greenberg22}, building on \citet{guerzoni03}, proposes that the scalar presupposition of \textit{only} is ``superlative'' so that $p$ is the weakest alternative in the context of focus alternatives, $C$. She also discusses how a general notion of stronger alternatives, symbolized as ‘$>$’, suffices, given that $C$ is similar to domain restriction on quantifiers and thus contains only contextually relevant alternatives, but at the same time the focus structure is constrained by the Question Under Discussion that is salient in the context (\citealt{rooth92}, \citealt{roberts96}, \citealt{beaverclark08}, \citealt{coppockbeaver10, coppockbeaver14}). Following \citet{guerzoni03}, \citet{greenberg22} treats \textit{only} and \textit{even} as antonyms, with mirroring superlative presuppositions, but adds special constraints on $C$ with respect to the QUD in order to account for the mirror image felicity patterns with \textit{only} and \textit{even}. She suggests that the superlative formulation of the scalar presupposition can be used to capture the differences in the types of scales with different operators, which is a question for future research. In this chapter, I focus on the common core of the exclusives, therefore I simply adopt the idea of stronger/weaker alternatives.

In (\ref{ex:tom:75}) above, I proposed a single lexical entry for \textit{tylko/zaledwie} because \textit{zaledwie} does not block \textit{tylko} from pragmatic scales and its use with numeral scales seems to only be dispreferred, as was shown in example \REF{ex:tom:37} in \sectref{sec:tom:4.2}. However, the approximative and counterfactual uses of \textit{(za)ledwie} require an extra presupposition: the lexical entry for \textit{(za)ledwie}\textsubscript{\textsc{approx}} makes it salient that \textit{not p} could happen, (\ref{ex:tom:76c}). I observed that those uses require obligatory wide focus (VP or sentence focus), which I take to correlate with the fact that $\neg p$ %\textit{not p} 
is a salient alternative. Given that $C$ is determined also by the QUD, if $\neg p$ %\textit{not} \textit{p} 
is salient, the context disambiguates between the lexical entries in \REF{ex:tom:75} and \REF{ex:tom:76}.

% \ea\label{ex:tom:76}
% \sib{\textit{(za)ledwie}$_{APPROX}$}$^w (C) (p)$ is defined iff\\
% (i) $\exists q \in C [ q \ge p \wedge q(w) = 1]$\\\hfill \textit{Factivity: an alternative at least as strong as $p$ is true}\\
% and\\
% (ii) $\forall q \in C [q \ne p \rightarrow q > p]$\hfill \textit{Scalarity: $p$ is the weakest alternative}\\
% and\\
% (iii) $\neg p > q$\hfill \textit{Scalarity: $p$ is weaker than $\neg p$}\\
% if defined then \sib{\textit{(za)ledwie}$_{APPROX}$}$^w (C) (p) = \forall q[(q \in C \wedge q \ne p \wedge q > p) \rightarrow q(w) = 0]$
% \z 

\ea\label{ex:tom:76}
\sib{\textit{(za)ledwie}\textsubscript{\textsc{approx}}}$^w(C)(p)$ is defined iff
\ea $\exists q\in C [q\geq p\wedge q(w) = 1]$\label{ex:tom:76a}\\{[Factivity: an alternative at least as strong as $p$ is true]}\smallskip\\   
and
\ex $\forall q\in C [q\neq p → q > p]$\label{ex:tom:76b}\\{[Scalarity: $p$ is the weakest alternative]}\smallskip\\
and
\ex $¬p > p$\label{ex:tom:76c}\\
{[Scalarity: $p$ is weaker than $¬p$]}\smallskip
\z
if defined then\\
\sib{\textit{(za)ledwie}\textsubscript{\textsc{approx}}}$^w(C)(p) =\forall q[(q\in C\wedge q\neq p\wedge q > p) → q (w) = 0]$ 
\z 

\noindent A separate entry for \textit{dopiero} is needed. This is the only exclusive that is limited to temporal scales and blocks \textit{tylko} and \textit{zaledwie} from temporal uses (in contrast to \textit{only} in English). In \sectref{sec:tom:4.3} we showed that in contrast to \textit{tylko}/\textit{zaledwie}, whose scalarity is pragmatically determined, \textit{dopiero} requires a temporal scale. We can build this into its semantics by keeping exactly the same meaning components as for \textit{tylko}/\textit{zaledwie} in \REF{ex:tom:75} but adding a presupposition specifically about a temporal scale such that the time of \textit{p} is more noteworthy than the times of the alternatives, either because it was early or because it was late. The exclusive assertion stays the same, because as I noted in \sectref{sec:tom:4.3}, \textit{dopiero} excludes stronger alternatives on the contextual scale, not necessarily later ones (e.g., buying more than 2 tomatoes, \REF{ex:tom:46}, vs. finishing earlier than in 2 months, \REF{ex:tom:62a}).\footnote{My intuition is that the exclusion of weaker alternatives should account for why the time of the event in the prejacent is more noteworthy than that of a stronger alternative and then the extra presupposition in (\ref{ex:tom:77c}) should be superfluous. This remains a topic for future work.}

% \ea\label{ex:tom:77} 
% \sib{\textit{dopiero}}$^w (C)(p)$ is defined iff\\
% (i) $\exists q \in C [q \ge p \wedge q(w) = 1]$\\\hfill \textit{Factivity: an alternative at least as strong as $p$ is true}\\
% and\\
% (ii) $\forall q \in C [q \ne p \rightarrow q > p]$\hfill \textit{Scalarity: $p$ is the weakest alternative}\\
% and\\ 
% (iii) $\exists t_1 \in T [p = \lambda w [ t_1\; \text{in}\; w]] \wedge \forall q \in C [\exists t_2 \in T [q = \lambda w [t_2\; \text{in}\; w]] \wedge q \ne p \wedge t_1 >_{NOTEWORTHY^w} t_2]$\hfill \textit{Temporality: Time $t$ of $p$ is the most noteworthy}\\
% if defined then \sib{\textit{dopiero}}$^w (C)(p) = \forall q [(q \in C \wedge q \ne p \wedge q > p) \rightarrow q(w) = 0$
% \z 

\ea\label{ex:tom:77} \sib{\textit{dopiero}}$^w(C)(p)$ is defined iff
\ea $\exists q \in C [q\geq p \wedge q(w) = 1]$\label{ex:tom:77a}\\{[Factivity: an alternative at least as strong as $p$ is true]}\smallskip\\
and
\ex $\forall q\in C[q\neq p → q > p]$\label{ex:tom:77b}\\
{[Scalarity: $p$ is the weakest alternative]}\smallskip\\
and
\ex $\exists t_1\in T[p = λw[t_1 \text{ in}\ w]]\wedge \forall q\in C[\exists t_2\in T[q = λw[t_2 \text{ in}\ w]]\wedge q \neq p\wedge t_1 >_{\textsc{noteworthy}^w} t_2]$\label{ex:tom:77c}\\
{[Temporality: Time $t$ of $p$ is the most noteworthy.]}\smallskip
\z
if defined then\\
\sib{\textit{dopiero}}$^w(C)(p)  =\forall q[(q\in C\wedge q\neq p\wedge q > p) → q (w) = 0]$ 
\z 

\noindent As established in \sectref{sec:tom:4.5}, \textit{a\.{z}} is the scalar opposite of \textit{tylko}/\textit{zaledwie}, thus the scale is reverted in the asserted and the presupposed components: \REF{ex:tom:78} vs. \REF{ex:tom:75}. The factive presupposition, (\ref{ex:tom:78a}), states that the prejacent or an alternative at most as strong is true. The scalar presupposition, \REF{ex:tom:78} (ii), is that the prejacent is the highest on the scale. Exclusivity but also scalarity, ‘no lower alternative’, is asserted in the last line of \REF{ex:tom:78}.

% \ea\label{ex:tom:78}
% \sib{\textit{aż}}$^w (C)(p)$ is defined iff\\
% (i) $\exists q \in C [q \le p \wedge q(w) = 1]$\\\hfill \textit{Factivity: an alternative at most as strong as $p$ is true}\\
% and\\
% (ii) $\forall q \in C [q \ne p \rightarrow p > q]$\hfill \textit{Scalarity: $p$ is the strongest alternative}\\
% if defined \sib{\textit{aż}}$^w (C)(p) = \forall q [(q \in C \wedge q \ne p \wedge p > q) \rightarrow q(w) = 0]$
% \z 

\ea\label{ex:tom:78}
\sib{\textit{aż}}$^w(C)(p)$ is defined iff
\ea $\exists q\in C [q\leq p\wedge q(w) = 1]$\label{ex:tom:78a}\\
{[Factivity: an alternative at most as strong as $p$ is true]}\smallskip\\
and
\ex $\forall q\in C [q\neq p → p > q]$\label{ex:tom:78b}\\
{[Scalarity: $p$ is the strongest alternative]}\smallskip    
\z
if defined then\\
\sib{\textit{aż}}$^w(C)(p) =\forall q[(q\in C\wedge q\neq p\wedge p > q ) → q (w) = 0]$ 
\z 

\noindent Compare now how \textit{aż} differs from the additive scalar adverb \textit{even}. The prejacent presupposition in (\ref{ex:tom:79a}) expresses additivity, and the scalar presupposition involves a scale of likelihood, (\ref{ex:tom:79b}) \citep{guerzoni03}. There is no contribution to the assertion of the sentence.

% \ea\label{ex:tom:79}
% \sib{\textit{even}}$^w (C)(p)$ is defined iff\\
% (i) $\exists q \in C [q \ne p \wedge q(w) = 1]$ \hfill \textit{Factivity and additivity: an alternative $q$ is true}\\
% and\\
% (ii) $\forall q \in C [q \ne p \rightarrow q >_{LIKELY^w} p]$\hfill \textit{Scalarity: $p$ is the least likely alternative}\\
% %and\\
% if defined then \sib{\textit{even}}$^w (C)(p) = p(w)$
% \z 

\ea\label{ex:tom:79}
\sib{\textit{even}}$^w(C)(p)$ is defined iff 
\ea $\exists q\in C [q\neq p\wedge q(w)=1]$\label{ex:tom:79a}\\
{[Factivity and additivity:  an alternative $q$ is true]}\smallskip\\
and
\ex $\forall q\in C [q\neq p → q >_{\textsc{likely}^w} p]$\label{ex:tom:79b}\\
{[Scalarity: $p$ is the least likely alternative]}\smallskip
\z
if defined then\\
\sib{\textit{even}}$^w(C)(p) = p(w)$
\z 

\noindent \citet{greenberg22} observes that although \textit{even} and \textit{only} are scalar opposites, \textit{even} must place the prejacent high on the scale, “above the norm” (formulated as a universal scalar presupposition in \citealt{guerzoni03}: the least likely among all the alternatives, (\ref{ex:tom:79b})). However, \textit{only} “does not inherently express ‘a little’, but just needs to indicate ‘less’ than its alternative(s) in C” \citep[23]{greenberg22}. For this reason \textit{only} is felicitous in the scenario in \REF{ex:tom:Greenberg-52a}, but \textit{even} is not felicitous in \REF{ex:tom:Greenberg-52b}.

\ea\label{ex:tom:Greenberg-52} 
\ea\label{ex:tom:Greenberg-52a}
John is tall. He is 1.92m. Bill is shorter -- he is only 1.88m (but he is still tall). \hfill \cite[(52a) p. 21]{greenberg22}
\ex\label{ex:tom:Greenberg-52b}
John is short. He is 1.52m. Bill is taller. He is (\#even) 1.58m (but he is still short). \hfill \cite[(52b) p. 21]{greenberg22}
\z
\z 

\noindent In contrast to \textit{even}/\textit{nawet}, \textit{aż} is indeed felicitous in the Polish counterpart of Greenberg’s example in \REF{ex:tom:Greenberg-plb}, thanks to its contextually defined scale, as opposed to a likelihood scale. (As expected, \REF{ex:tom:Greenberg-pla} shows that \textit{tylko}/\textit{zaledwie} behave in the exact same way as \textit{only}.)

\ea\label{ex:tom:Greenberg-pl} 
\ea\label{ex:tom:Greenberg-pla}
\gll Jan jest wysoki. Ma 192 cm wzrostu. Bill jest niższy. Ma \minsp{\{} tylko / zaledwie\} 188 cm wzrostu (ale i tak jest wysoki.)\\
Jan is tall has 192 cm height Bill is shorter has {} only {} merely 188 cm height but even so is tall\\
\glt ‘Jan is tall. He is 192 cm. Bill is shorter -- he is only 188 cm (but he is still tall).’ 
\ex\label{ex:tom:Greenberg-plb}
\gll Jan jest niski. Ma 152 cm wzrostu. Bill jest wyższy. Ma \minsp{\{} aż / \minsp{\#} nawet\} 158 cm wzrostu (ale i tak jest niski.)\\
Jan is short has 152 cm height Bill is taller has {} \textsc{aż} {} {} even 158 cm height but even so is short \\
\glt ‘Jan is short. He is 152 cm. Bill is taller. He is \textit{aż}/(\#even) 158 cm (but he is still short).’ 
\z
\z 

\noindent Moreover, it can be demonstrated that \textit{aż} cannot use a likelihood scale that matches the wider discourse/world-knowledge. The scale with \textit{aż} must be constructed with reference to the lexical content of the associate, e.g., a non-pro\-fes\-sion\-al in \REF{ex:tom:80}. In \REF{ex:tom:80} \textit{even} adds the reading that hiring a non-professional is very unlikely (under normal circumstances). \textit{Aż} is infelicitous in this sentence because it requires that a non-professional is ranked high on the scale of people we want to hire.

\ea\label{ex:tom:80}
\gll Potrzebujemy sobowtóra głównego  aktora. Zatrudnimy \minsp{\{} nawet / \minsp{\#} aż\}   nieprofesjonalistę.\\
we.need body.double lead actor will.hire {} even {} {} \textsc{aż} non.professional\\
\glt `We need a body double of the lead actor. We will even hire a non-professional.’
\z

\noindent The factive presupposition of \textit{aż}, (\ref{ex:tom:78a}), is not additive like the presupposition of \textit{even}, (\ref{ex:tom:79a}). This means that \textit{aż} can occur in additive contexts, \REF{ex:tom:applesa}--\REF{ex:tom:applesb}, and it can even co-occur with \textit{nawet} `even', \REF{ex:tom:head}.\footnote{I thank the anynonymous reviewer for raising this issue.} 

\ea\label{ex:tom:apples} 
\ea\label{ex:tom:applesa} 
\gll Janek zjadł aż 5 jabłek.\\
Janek ate \textsc{aż} 5 apples\\
\glt `Janek ate even 5 apples.’
\ex\label{ex:tom:applesb} 
\gll Janek nie rozmawiał zaledwie z dziekanem, lecz aż z rektorem.\\
Janek not talked merely with dean but \textsc{aż} with chancellor\\
\glt `Janek did not talk merely to the dean but to the chancellor himself.’\\\hfill \hfill (\citealt{tomaszewicz12}: 325, ex. (7b))
\z 
\z 

\ea\label{ex:tom:head}  
\gll Maria rozmawiała z szefem zakładu, rozmawiała z dziekanem, rozmawiała nawet aż z rektorem.\\
Maria talked with head institute talked with dean talked even \textsc{aż} with chancellor\\
\glt `Maria talked to the institute head, she talked to the dean, she talked even to the chancellor himself.’
\z

\noindent The emotive factive test in \REF{ex:tom:head-emotive} shows that with \textit{nawet} the assertion is the prejacent, but with \textit{aż} it is the excluded component, namely, Maria’s happiness in \REF{ex:tom:head-emotiveb} is about the fact that she got to talk to somebody high on the scale of important people, not about the fact that she talked to these particular people (the same effect was shown in \REF{ex:tom:64}). 

\ea\label{ex:tom:head-emotive}  
\gll Maria rozmawiała z szefem zakładu, rozmawiała z dziekanem, …\\
Maria talked with head institute talked with dean\\
\glt ‘Maria talked to the institute head, she talked to the dean, …'
\ea\label{ex:tom:head-emotivea}  
\gll i jest zadowolona, bo rozmawiała   nawet z rektorem.\\
and is happy because talked even with chancellor\\
\glt ‘and she's happy because she talked even to the chancellor.'\\
\textit{Implied:} ‘Maria is happy because she talked to the chancellor.’
\ex\label{ex:tom:head-emotiveb}  
\gll i jest zadowolona, bo rozmawiała aż z rektorem.\\
and is happy because talked \textsc{aż} with chancellor\\
\glt ‘and she's happy because she talked to the chancellor himself.'\\
\textit{Implied:} ‘Maria is happy because she talked to somebody no less important than the chancellor.’
\z 
\z 

\noindent This also indicates that a shift in the QUD (from `Who did Mary talk to?' to `How important a person did Mary talk to?') changes the set of relevant alternatives in $C$, pruning some of those included in the prior discourse as shown in \citet{greenberg22} for both \textit{even} and \textit{only}. In \REF{ex:tom:applesa} \textit{aż} is felicitous although it is true that Janek ate 4/3/2/1 apples.

The formal mechanism of focus association derives the set of alternatives in $C$ (so that the size and location of focus determine the contents of the set) and this set can be further restricted by the current QUD. Below, I illustrate this for \textit{aż}; the mechanism works exactly the same for any focus-associating item taking sentential scope.

% \ea\label{ex:tom:81}
% Condition on focus association \citep{rooth92}\\
% $C \subseteq$ \sib{$\alpha$}$^f$, or $C \subseteq \cup$ \sib{$\alpha$}$^f$, where $C$ is the restrictor of a quantificational adverb and $alpha$ the sister to $\sim$
% \z 

\ea\label{ex:tom:81}
Condition on focus association \citep{rooth92}\smallskip\\
$C\subseteq\text{\sib{α}}^f$, or $C\subseteq\bigcup\text{\sib{α}}^f$\smallskip\\
where \textit{C} is the restrictor of a quantificational adverb and α the sister to $\sim$
\z 

\noindent In \REF{ex:tom:82} I present the derivation for association with the focus ‘Warsaw’. The meaning of the sentence in \REF{ex:tom:82c} results from the specification of the lexical entry in \REF{ex:tom:78}. Note how it does not tell us what the weaker alternatives are. This information is derived from focus association: the focus operator $\sim$ that interprets the feature F on the syntactic constituent is attached at the sentential level, \REF{ex:tom:82b}, and thus its restrictor variable $S$ is a subset of the focus alternative value of the TP. The focal presupposition is the set of alternative propositions \{`that Mary drove to Łódź', `that Mary drove to Grójec', `that Mary drove to Berlin', \dots\} in \REF{ex:tom:82d}. The restrictor of \textit{aż}, $C$, is a subset of that set satisfying the condition on focus association in \REF{ex:tom:81}.  Because $q \in C$, the excluded alternatives are a subset of the set in \REF{ex:tom:82d}. This subset does not include stronger alternatives such as `that Mary drove to Berlin' because Berlin is further than Warsaw in the context where Mary is driving from Bełchatów (as seen in Figure \ref{fig:tom:warsaw}). Obviously, on her way to Warsaw Mary drove through other places, but the QUD is `How significantly far did Mary drive?' so the alternatives, `Mary drove to Łódź', `Mary drove to Grójec' are true, but do not answer the QUD because those places are not significantly far from Bełchatów. 

\ea\label{ex:tom:82}
% \ea\label{ex:tom:82a} 
\gll Mary pojechała aż do \minsp{[} Warszawy]\textsubscript{F}.   \\
Mary drove \textsc{aż} to {} Warsaw\\
\glt ‘Mary drove all the way to Warsaw.'
% \ex\label{ex:tom:82b}     
% {}[\textsc{aż} $C$] [\textsubscript{TP\textsubscript{2}} [$\sim$ $S$] [\textsubscript{TP\textsubscript{1}} Mary drove to [Warsaw]\textsubscript{F}]] 
% \ex\label{ex:tom:82c}     
% \sib{\REF{ex:tom:82a}} $= \lambda w.\forall q [(q \in C \wedge q \ne$ \sib{Mary drove to Warsaw} $\wedge\; q <$ \sib{Mary drove to Warsaw} $\rightarrow \neg q(w)]$
% \ex\label{ex:tom:82d}     
% $S \subseteq$ \sib{TP\textsubscript{1}}$^f = \{q: \exists x [ q = \lambda w [\text{drive-to}(m, x)(\text{in}\; w)]]\}$
% \ex\label{ex:tom:82e}
% $C \subseteq$ \sib{TP\textsubscript{1}}$^f$ \hfill (\textit{focus association}, \REF{ex:tom:81})
% \z 
\ea\label{ex:tom:82b}     
{}[\textsc{aż} \textit{C}] [\textsubscript{TP$_2$} [{$\sim$} \textit{S}] [\textsubscript{TP$_1$} Mary drove to [Warsaw]\textsubscript{F}]] 
\ex\label{ex:tom:82c}     
\sib{\REF{ex:tom:82}}${}= λw.\forall q[(q\in C\wedge q\neq{}$\sib{\text{Mary drove to Warsaw}}${}\wedge q <{}$\sib{\text{Mary drove to Warsaw}}${} → ¬ q (w)]$
\ex\label{ex:tom:82d}     
$S\subseteq{}$\sib{TP\textsubscript{1}}$^f = \{q: \exists x [q = λw [\textsc{drive to}(\textsc{Mary},x)\text{ in }w)]]\}$
\ex\label{ex:tom:82e}
$C\subseteq{}$\sib{TP\textsubscript{1}}$^f$\hfill [focus association; see \REF{ex:tom:81}]
\z
\z 

\noindent When \textit{aż} associates with VP focus, the set C is a subset of the set of different things that Mary did, \REF{ex:tom:83d}.

\ea\label{ex:tom:83}
% \ea\label{ex:tom:83a} 
\gll Mary aż \minsp{[} pojechała do Warszawy]\textsubscript{F}.  \\
Mary \textsc{aż} {} drove to Warsaw\\
\glt ‘Mary even drove to Warsaw (of all things she could do).'
% \ex\label{ex:tom:83b}      
% {}[\textsc{aż} $C$] [\textsubscript{TP\textsubscript{2}} [$\sim$ $S$] [\textsubscript{TP\textsubscript{1}} Mary [drove to Warsaw]\textsubscript{F}]]
% \ex\label{ex:tom:83c}      
% \sib{\REF{ex:tom:83a}} $= \lambda w.\forall q [(q \in C \wedge q \ne$ \sib{Mary drove to Warsaw} $\wedge\; q <$ \sib{Mary drove to Warsaw} $\rightarrow \neg q(w)]$
% \ex\label{ex:tom:83d} 
% $S \subseteq$ \sib{TP\textsubscript{1}}$^f = \{q: \exists f \exists x [ q = \lambda w [f(m, x)(\text{in}\; w)]]\}$
% \ex\label{ex:tom:83e} 
% $C \subseteq$ \sib{TP\textsubscript{1}}$^f$ \hfill (\textit{focus association}, \REF{ex:tom:81})
\ea\label{ex:tom:83b}      
{}[\textsc{aż} \textit{C}] [\textsubscript{TP$_2$} [{$\sim$} \textit{S}] [\textsubscript{TP$_1$} Mary [drove to Warsaw]\textsubscript{F}]]
\ex\label{ex:tom:83c}      
\sib{\REF{ex:tom:83}}${}= λw.\forall q[(q \in C\wedge q\neq{}$\sib{Mary drove to Warsaw}${}\wedge q <{}$\sib{Mary drove to Warsaw})${}→ q(w)]$
\ex\label{ex:tom:83d} 
$S\subseteq{}$\sib{TP\textsubscript{1}}$^f = \{q: \exists f \exists x [q = λw [f(\textsc{Mary}, x)$ in $w]]\}$  
\ex\label{ex:tom:83e} 
$C\subseteq{}$\sib{TP\textsubscript{1}}$^f$\hfill [focus association; see \REF{ex:tom:81}]
\z 
\z 

\noindent When \textit{aż} associates with sentence focus, the alternatives are all alternative things that could happen, \REF{ex:tom:84d}.

\ea\label{ex:tom:84}
% \ea\label{ex:tom:84a} 
\gll Aż \minsp{[} Mary pojechała do Warszawy]\textsubscript{F}.   \\
\textsc{aż} {} Mary drove to Warsaw\\
\glt ‘Of all things that could happen, Mary drove to Warsaw.'
% \ex\label{ex:tom:84b}      
% {}[\textsc{aż} $C$] [\textsubscript{TP\textsubscript{2}} [$\sim$ $S$] [\textsubscript{TP\textsubscript{1}} [Mary drove to Warsaw]\textsubscript{F}]]
% \ex\label{ex:tom:84c} 
% \sib{\REF{ex:tom:84a}} $= \lambda w.\forall q [(q \in C \wedge q \ne$ \sib{Mary drove to Warsaw} $\wedge\; q <$ \sib{Mary drove to Warsaw} $\rightarrow \neg q(w)]$
% \ex\label{ex:tom:84d}      
% $S \subseteq$ \sib{TP\textsubscript{1}}$^f = \{q \in D_{\langle s, t \rangle}\}$
% \ex\label{ex:tom:84e} 
% $C \subseteq$ \sib{TP\textsubscript{1}}$^f$ \hfill (\textit{focus association}, \REF{ex:tom:81})
\ea\label{ex:tom:84b}      
[\textsc{aż} \textit{C}] [\textsubscript{TP$_2$} [$\sim$ \textit{S}] [\textsubscript{TP$_1$} [Mary drove to Warsaw]\textsubscript{F}]]
\ex\label{ex:tom:84c}      
\sib{\REF{ex:tom:84}}${}= λw.\forall q[(q\in C\wedge q\neq{}$\sib{Mary drove to Warsaw}${}\wedge q <{}$ \sib{Mary drove to Warsaw})${}→ ¬ q(w)]$
\ex\label{ex:tom:84d}      
$S\subseteq{}$ \sib{TP\textsubscript{1}}$^f=\{q\in D_{\stb{s,t}}\}$
\ex\label{ex:tom:84e} 
$C\subseteq{}$\sib{TP\textsubscript{1}}$^f$\hfill [focus association; see \REF{ex:tom:81}]
\z 
\z 

\noindent Since \textit{aż} may but does not have to have a temporal reading like \textit{dopiero}, I propose that the temporal interpretation is determined by pragmatics and syntax. Namely, \textsc{temporal} \textit{aż} functions as a complementizer in Polish so it precedes a clause and the focus association mechanism derives a set of alternative propositions which are ordered with respect to time. As noted in footnote \ref{fn:tom:former4} the semantics and syntax alone cannot be the reason why \textit{aż} can function as a complementizer but the other exclusive particles cannot. I propose that when \textit{aż} associates with the whole sentence, as in \REF{ex:tom:85c}, the set of focus alternatives is the set of answers to the QUD, and pragmatics determines that the temporal ordering is the most relevant. What possibly contributes to this interpretation is the fact that without context a sort of default ordering for multiple events is a temporal ordering. In other words, I assume that, unlike for \textit{dopiero}, for \textsc{temporal} \textit{aż} no extra presupposition is necessary in the lexical entry. The set $C$ in \REF{ex:tom:85e} evokes multiple events and times for the set of propositions derived from wide focus. So what ends up being excluded are alternative propositions describing different events that happened at different times. The scalar presupposition says that those times are ordered on a scale. The truth conditions say that no weaker proposition, i.e. nothing that happened earlier, is true.\footnote{There is an interesting contrast in Polish between two types of \textit{until}-clauses: \textit{aż} and \textit{dopóki}. In contrast to \textit{aż}, \REF{ex:tom:85} repated as \REF{ex:tom:foota} below, \textit{dopóki} requires negation for an identical interpretation, \REF{ex:tom:footb} and my intuition is that this is because no exclusivity is built into its semantics. 
\ea\label{ex:tom:foot}
\ea\label{ex:tom:foota} 
\gll Zostaniemy tu aż zajdzie słońce. = \REF{ex:tom:85}\\
will.stay here \textsc{aż} goes.down sun\\
\glt `We will stay here until the sun goes down.’
\ex\label{ex:tom:footb} 
\gll Zostaniemy tu dopóki nie zajdzie słońce. \\
will.stay here \textsc{dopóki} not goes.down sun\\
\glt `We will stay here until the sun goes down.’
\z 
\z 
Since \textit{aż} excludes lower alternatives on the time scale, we obtain the time point where all alternatives to the sun going down are no longer true. But if \textit{dopóki} works like ‘as long as’, negation is necessary to compositionally derive the endpoint. For extensive work on the Czech counterpart of \textit{dopóki}, \textit{dokud}, see \citet{docekal12a,docekal15}, \citet{iatridouzeijlstra21}.
}

\ea\label{ex:tom:85}
% \ea\label{ex:tom:85a} 
\gll Zostaniemy tu aż zajdzie słońce.\\
will.stay here \textsc{aż} goes.down sun\\
\glt `We will stay here until the sun goes down.’
% \ex\label{ex:tom:85b}
% Aż [sun goes down]\textsubscript{F}.
% \ex\label{ex:tom:85c}  
% {}[\textsc{aż} $C$] [\textsubscript{TP\textsubscript{2}} [$\sim$ $S$] [\textsubscript{TP\textsubscript{1}} [sun goes down]\textsubscript{F}]]
% \ex\label{ex:tom:85d}     
% \sib{\REF{ex:tom:85}} $= \lambda w.\forall q [(q \in C \wedge q \ne$ \sib{sun goes down} $\wedge\; q <$ \sib{sun goes down} $\rightarrow \neg q(w)]$
% \ex\label{ex:tom:85e}      
% $S \subseteq$ \sib{TP\textsubscript{1}}$^f = \{q: \exists t \in T [q = \lambda w [ t\; \text{in}\, w]]\}$
% \ex\label{ex:tom:85f} 
% $C \subseteq$ \sib{TP\textsubscript{1}}$^f$ \hfill (\textit{focus association}, \REF{ex:tom:81})
\ea\label{ex:tom:85b}      
Aż [sun goes down]\textsubscript{F}.
\ex\label{ex:tom:85c}     
{}[\textsc{aż} \textit{C}] [\textsubscript{TP2} [$\sim$ \textit{S}] [\textsubscript{TP1} [sun goes down]\textsubscript{F}]]
\ex\label{ex:tom:85d}     
\sib{\REF{ex:tom:85}} = $λw.\forall q [(q \in C\wedge q\neq{}$\sib{sun goes down}${}\wedge q <{}$\sib{sun goes down}$) → ¬q(w))]$
\ex\label{ex:tom:85e}      
$S\subseteq{}$\sib{TP\textsubscript{1}}$^f = \{q:\exists t\in T [q = λw [t$ in $w]]\}$
\ex\label{ex:tom:85f} 
$C\subseteq{}$\sib{TP\textsubscript{1}}$^f$\hfill [focus association; see \REF{ex:tom:81}]
\z 
\z 

\noindent Finally, let us see why the four exclusives are interchangeable in the sentence in \REF{ex:tom:1}, repeated below in \REF{ex:tom:86}. The logical syntax is given in \REF{ex:tom:87a}. The focus is on the numeral and the exclusive takes sentential scope. This derives the set of alternatives in \REF{ex:tom:87c}. With \textit{tylko/zaledwie/dopiero} the truth conditions say that all of the stronger alternatives are excluded, e.g. ‘that the team achieved 1\textsuperscript{st} place, that the team achieved 15\textsuperscript{th} place, …’, \REF{ex:tom:87b} and \REF{ex:tom:88b}. The scalar presupposition is that the prejacent ‘that the team achieved 30\textsuperscript{th} place’ is the weakest alternative. 

\ea\label{ex:tom:86}
\gll Zespół był słaby i zajął \minsp{\{} tylko / zaledwie / dopiero / aż\}   \minsp{[} 30]\textsubscript{F} miejsce.\\
team was poor and achieved {} only {} merely {} only\textsubscript{\textsc{temp}} {} \textsc{aż} {} 30\textsuperscript{th} place\\\hfill (cf. \REF{ex:tom:1})
\glt `The team performed poorly and earned only 30\textsuperscript{th} place.’
\z 

% \ea\label{ex:tom:87}
% \ea\label{ex:tom:87a} 
% {}[\textsc{tylko/zaledwie} $C$] [\textsubscript{TP\textsubscript{2}} [$\sim$ $S$] [\textsubscript{TP\textsubscript{1}} the team earned [30\textsuperscript{th}]\textsubscript{F} place]]
% \ex\label{ex:tom:87b} 
% \sib{\REF{ex:tom:87a}} $= \lambda w.\forall q [(q \in C \wedge q \ne$ \sib{the team earned 30\textsuperscript{th} place} $\wedge\; q >$ \sib{the team earned 30\textsuperscript{th} place} $\rightarrow \neg q(w)]$
% \ex\label{ex:tom:87c}      
% $S \subseteq$ \sib{TP\textsubscript{1}}$^f = \{q: \exists n \in \mathbb{N} [q = \lambda w. \text{The team earned}\; n^{th}\; \text{place in}\; w]\}$
% \ex\label{ex:tom:87d} 
% $C \subseteq$ \sib{TP\textsubscript{1}}$^f$ \hfill (\textit{focus association}, \REF{ex:tom:81})
% \z 
% \z 

\ea\label{ex:tom:87}
\ea\label{ex:tom:87a} 
{}[\textsc{tylko/zaledwie} \textit{C}] [\textsubscript{TP$_2$} [$\sim$ \textit{S}] [\textsubscript{TP$_1$} the team earned [30\textsuperscript{th}]\textsubscript{F} place]]
\ex\label{ex:tom:87b}      
\sib{\REF{ex:tom:87a}}${} = λw.\forall q [(q\in C\wedge q\neq{}$\sib{the team earned 30\textsuperscript{th} place}${}\wedge q >{}$\sib{the team earned 30\textsuperscript{th} place}$) → ¬q(w)]$
\ex\label{ex:tom:87c}      
$S\subseteq{}$\sib{TP\textsubscript{1}}$^f = \{q:\exists n\in ℕ [q = λw.{}$\textsc{the team earned} \textit{n}\textsuperscript{th} \textsc{place} in $w]\}$
\ex\label{ex:tom:87d} 
$C\subseteq{}$\sib{TP\textsubscript{1}}$^f$\hfill [focus association; see \REF{ex:tom:81}]
\z\z 

\noindent With \textit{dopiero} the excluded alternatives are further restricted by the presupposition about noteworthiness of alternative times, \REF{ex:tom:88d}, e.g., ‘that the team achieved 1\textsuperscript{st} place at time $t_1$, that the team achieved 15\textsuperscript{th} place at time $t_{15}$, …’. So although the inherent ordering between numerals is such that 30 $>$ 15, the context defines the scale where achieving 15\textsuperscript{th} is a stronger alternative than achieving 30\textsuperscript{th} place. 


% \ea\label{ex:tom:88}
% \ea\label{ex:tom:88a} 
% {}[\textsc{dopiero} $C$] [\textsubscript{TP\textsubscript{2}} [$\sim$ $S$] [\textsubscript{TP\textsubscript{1}} the team earned [30\textsuperscript{th}]\textsubscript{F} place]]
% \ex\label{ex:tom:88b}     
% \sib{\REF{ex:tom:88a}} $= \lambda w.\forall q [(q \in C \wedge q \ne$ \sib{the team earned 30\textsuperscript{th} place} $\wedge\; q >$ \sib{the team earned 30\textsuperscript{th} place} $\rightarrow \neg q(w)]$
% \ex\label{ex:tom:88c}      
% $S \subseteq$ \sib{TP\textsubscript{1}}$^f = \{q: \exists n \in \mathbb{N} [q = \lambda w. \text{The team earned}\; n^{th}\; \text{place in}\; w]\}$
% \ex\label{ex:tom:88d} 
% $C \subseteq$ \sib{TP\textsubscript{1}}$^f$ \hfill (\textit{focus association}, \REF{ex:tom:81})
% \ex\label{ex:tom:88e}
% $\exists t_1 \in T [p = \lambda w.$ The team earned 30\textsuperscript{th} place at $t_1$ in $w$] $\wedge \forall q \in C [\exists t_2 \in T\; \exists n \in \mathbb{N} [ q = \lambda w.$ The team earned $n^{th}$ place at $t_2$ in $w]%\}
% ]$ $\wedge q \ne p \wedge t_1 >_{NOTWORTHY^w} t_2]$ \\\hfill \textit{(presupposition in (\ref{ex:tom:77c}) restricting the set in \REF{ex:tom:88d})}
% %${\exists}$\textit{t}\textsubscript{1} \textrm{${\in}$ T} [p = λ\textit{w}.  The team earned \textit{30\textsuperscript{th}} place at \textit{t}\textsubscript{1} in \textit{w} ] \textrm{${\wedge}$ ${\forall}$}q \textrm{${\in}$}C [${\exists}$\textit{t}\textsubscript{2} \textrm{${\in}$ T} ${\exists}$\textit{n} \textrm{${\in}$} \textrm{$\mathbb{N}$} [\textit{q} = λ\textit{w}. The team earned \textit{n\textsuperscript{th}} place at \textit{t}\textsubscript{2} in \textit{w}]\}] \textrm{${\wedge}$} q ${\neq}$ p  \textrm{${\wedge}$} \textit{t}\textsubscript{1} >\textsubscript{NOTEWORTHY} \textsuperscript{W} \textsubscript{} \textit{t}\textsubscript{2} ] (\textit{presupposition in \REF{ex:tom:77c} restricting the set in \REF{ex:tom:88d}})
% \z 
% \z 

\ea\label{ex:tom:88}
\ea\label{ex:tom:88a} 
{}[\textsc{dopiero} \textit{C}] [\textsubscript{TP$_2$} [$\sim$ \textit{S}] [\textsubscript{TP$_1$} the team earned [30\textsuperscript{th}]\textsubscript{F} place]]
\ex\label{ex:tom:88b}     
\sib{\REF{ex:tom:88a}}${} = λw.\forall q [(q\in C\wedge q\neq{}$\sib{the team earned 30\textsuperscript{th} place}${}\wedge q >{}$\sib{the team earned 30\textsuperscript{th} place}$) → ¬q(w)]$
\ex\label{ex:tom:88c}      
$S\subseteq{}$\sib{TP\textsubscript{1}}$^f = \{q:\exists n\in ℕ [q = λw.{}$\textsc{the team earned} \textit{n}\textsuperscript{th} \textsc{place} in $w]\}$
\ex\label{ex:tom:88d} 
$C\subseteq{}$\sib{TP\textsubscript{1}}$^f$\hfill [focus association; see \REF{ex:tom:81}]
\ex\label{ex:tom:88e}
$\exists t_1\in T [p = λw.{}$\textsc{the team earned} 30\textsuperscript{th} 
\textsc{place} at $t_1$ in $w]\wedge{}$ $\forall q\in C [\exists t_2\in T\;\exists n\in ℕ [q = λw.{}$ \textsc{the team earned} \textit{n}\textsuperscript{th} \textsc{place} at $t_2$ in $w]]\wedge q\neq p\wedge t_1>_{\textsc{noteworthy}^w}t_2]$\\\hfill [presupposition in \REF{ex:tom:77c} restricting the set in \REF{ex:tom:88d}]
\z \z 

\noindent With \textit{aż} the scalar presupposition is different: the prejacent is the strongest alternative and the ordering in $C$ is  ‘that the team achieved 30\textsuperscript{th} place’ $>$ ‘that the team achieved 15\textsuperscript{th} place’. The wider context makes it clear that achieving 30\textsuperscript{th} place is ‘less desirable’, ‘less than expected’ than achieving a place lower on the numeral scale.

% \ea\label{ex:tom:89}
% \ea\label{ex:tom:89a} 
% {}[\textsc{aż} $C$] [\textsubscript{TP\textsubscript{2}} [$\sim$ $S$] [\textsubscript{TP\textsubscript{1}} the team earned [30\textsuperscript{th}]\textsubscript{F} place]]
% \ex\label{ex:tom:89b}
% \sib{\REF{ex:tom:89a}} $= \lambda w.\forall q [(q \in C \wedge q \ne$ \sib{the team earned 30\textsuperscript{th} place} $\wedge\; q >$ \sib{the team earned 30\textsuperscript{th} place} $\rightarrow \neg q(w)]$
% \ex\label{ex:tom:89c}      
% $S \subseteq$ \sib{TP\textsubscript{1}}$^f = \{q: \exists n \in \mathbb{N} [q = \lambda w. \text{The team earned}\; n^{th}\; \text{place in}\; w]\}$
% \ex\label{ex:tom:89d} 
% $C \subseteq$ \sib{TP\textsubscript{1}}$^f$ \hfill (\textit{focus association}, \REF{ex:tom:81})
% \z 
% \z 

\ea\label{ex:tom:89}
\ea\label{ex:tom:89a} 
{}[\textsc{aż} \textit{C}] [\textsubscript{TP$_2$} [$\sim$ \textit{S}] [\textsubscript{TP$_1$} the team earned [30\textsuperscript{th}]\textsubscript{F} place]]
\ex\label{ex:tom:89b}
\sib{\REF{ex:tom:89a}}${} = λw.\forall q [(q\in C\wedge q\neq{}$\sib{the team earned 30\textsuperscript{th} place}${}\wedge q <{}$\sib{the team earned 30\textsuperscript{th} place}$) → ¬q(w)]$
\ex\label{ex:tom:89c}      
$S\subseteq{}$\sib{TP\textsubscript{1}}$^f = \{q:\exists n\in ℕ [q = λw.{}$\textsc{the team earned} \textit{n}\textsuperscript{th} \textsc{place} in $w]\}$
\ex\label{ex:tom:89d} 
$C\subseteq{}$\sib{TP\textsubscript{1}}$^f$\hfill [focus association; see \REF{ex:tom:81}]
\z\z 

\noindent Thus, the example in \REF{ex:tom:1} is a special case where the truth conditions end up the same with opposite scalar orderings of alternatives, thanks to the context.

\section{Family members in other Slavic languages}
\label{sec:tom:6}

In \sectref{sec:tom:3}, I discussed the different syntactic configurations for focus association and noted that not all of them may be available in every language; e.g., in English \textit{only} in the clause-initial position cannot associate with the whole clause as it is possible in Polish or German. Now that I identified the minimal semantic differences between the four exclusives in Polish, we can expect that those same parameters will have different settings cross-linguistically and not all of the options may be available or, conversely, there could be more options for variation than in Polish. 

For instance, the limits on scale dimension that are lacking with %the 
English \textit{only}, are also found with %the 
German \textit{erst} \citep{koenig91}, a temporal exclusive like \textit{dopiero}. Several other Slavic languages have a counterpart of \textit{aż}, \textit{až} in Czech, Slovak and Russian, and \textit{čak} in Bulgarian, Serbian and other South Slavic languages (\citealt{tomaszewicz13,tomaszewicz13a}), but they seem to have various requirements on the type of scale. While in Russian \textit{až} is highly expressive and its use is considered colloquial, at first sight, the same range of contexts I have considered here for Polish seems to be available except for the two scope options with negation (Katja Jasinskaja, p.c.). In Czech and Slovak, on the other hand, there seems to be a strong requirement for temporal or distance scales, specifically with DP focus. Thus \REF{ex:tom:90} in Czech does not have the reading that the manager is a noteworthy person for Maria to talk to, but instead a temporal scale needs to be accommodated, as indicated by the translation (Radek Šimík, p.c.)\footnote{There may be variation among the speakers. As reported in \citealt{tomaszewicz13a} for the example in \REF{ex:tom:president} some of my informants preferred  the scalar interpretation in (a), rather than the temporal interpretation in (b). A dialectal difference is possible (Věra Dvořák, p.c.) – the speakers who preferred \REF{ex:tom:president}a are
from Moravia.

\ea\label{ex:tom:president}
\gll Hanka poprosila o pomoc až prezidenta. \hfill (Czech; \cite[(54)]{tomaszewicz13a})
\\ 
Hanka asked for help \textsc{až} president\\
\ea ‘It was nobody less important than the president who Hanka asked for help.’
\ex ‘It was no sooner than when talking to the president that Hanka asked for help.’\z 
\z
}. In Polish this reading is obtained with \textit{dopiero} in this sentence.

\ea\label{ex:tom:90}
\gll Maria mluvila až s manažerem. \\ 
Maria talked \textsc{až} with manager\\
\ea ‘Maria talked to the manager who appeared late and Maria talked to him but not the people who appeared before him.’
\ex \textit{Unavailable:} ‘Maria talked to somebody so important as the manager.’\hfill (Czech; Radek Šimík, p.c.)
\z 
\z 

\noindent With temporal scales \textit{až} in Czech has the same exclusive reading as %the 
Polish \textit{zaledwie}, \textit{dopiero} and \textit{aż} in the example we saw in \REF{ex:tom:32}, i.e., where it can be translated as \textit{only} in English but is incompatible with Polish \textit{tylko}. The sentence in \REF{ex:tom:91} is the counterpart of the Polish \REF{ex:tom:32}. Just like in Polish, the reading with \textit{až} is different from the reading with \textit{even}. With \textit{až} the sentence tells us that 6 o’clock is a noteworthy time to wake up, because it is late, i.e., ‘no sooner than’. With \textit{even} the meaning is ‘Dan woke up several times in addition to the unlikely time at 6 AM’.

\ea\label{ex:tom:91}
\gll Daneček se vzbudil až v 6 ráno.\\
Dan self woke.up \textsc{až} at 6 morning \\
\ea `Little Dan woke up only/no sooner than at 6 AM.’ 
\ex \textit{Unavailable:} `Little Dan woke up even at 6 AM.’\\\hfill (Czech; \cite[(6)]{tomaszewicz13a}) 
\z 
\z 

\noindent However, \textit{až} in Czech can also associate with a VP and evoke a pragmatic scale (\citealt{tomaszewicz13,tomaszewicz13a}). In \REF{ex:tom:92}
% \citep[from][]{capek91}
the alternatives are different things that could result from the boy’s surprise; there is no implication of a temporal ordering.

\ea\label{ex:tom:92}
\gll Chlapec v zeleném overalu se překvapením až [\textsubscript{VP} začervenal].\\
boy in green overalls self surprise \textsc{až} {} reddened\\% \citep{capek91}
\glt `The boy in a green overalls as much as blushed with surprise.’\\\hfill (Czech; K. Čapek: \textit{Od člověka k člověku III})
\z 

\noindent In Slovak \textit{až} can have temporal readings when modifying a DP with a temporal reading, such as ‘Friday’ in \REF{ex:tom:sl1a}--\REF{ex:tom:sl1b}, without the high-on-the-scale implication that is found in such cases in Polish, i.e., in both \REF{ex:tom:sl1a}--\REF{ex:tom:sl1b} the meaning is `until Friday' with the implication that Friday is included (`through Friday' in American English). Compare how the Polish counterpart of \REF{ex:tom:sl1a} contains a preposition \textit{do} ‘\textit{to’}, \REF{ex:tom:plfriday}, and in the presence of \textit{aż} (‘\textit{aż do Piątku’}) the meaning would be that Friday is late, which would be very odd in this context (we saw this effect in \REF{ex:tom:62c}). The temporal \textit{až} in Slovak is reminiscent of the complementizer \textit{aż} in Polish, in that a temporal endpoint is specified without evoking noteworthiness.

\ea\label{ex:tom:sl1}
\ea\label{ex:tom:sl1a} 
\gll V dňoch pondelok až piatok bude predajňa zatvorená.\\
in days Monday \textsc{až} Friday \textsc{aux} store closed\\
\glt `In days from Monday to Friday, the store will be closed.'
\ex\label{ex:tom:sl1b}            
\gll Keby budeš budúci týždeň na služobke? Utorok až piatok.\\
when \textsc{aux} next week on business.trip Tuesday \textsc{až} Friday\\
\glt `When will you be on a business trip next week? Tuesday through Friday.'\\\hfill (Slovak; Andrea Kellö Žačoková, Jana Bašnáková, p.c.)
\z 
\z

\ea\label{ex:tom:plfriday}
\gll W dniach od poniedziałku do piątku sklep będzie zamknięty.\\
in days from Monday to Friday store \textsc{aux} closed\\
\glt `In days from Monday to Friday, the store will be closed.'
\z

\noindent The following example, \REF{ex:tom:sl2}, shows that in Slovak \textit{až} can be used with intervals in general (not just with temporal endpoints as above) and again in Polish in such cases the preposition \textit{do} must be used without \textit{až}, \REF{ex:tom:pl400}, or else the sentence would carry the implication that 400 is significantly more than 384 pages.

\ea\label{ex:tom:sl2}
\gll Naša kniha bude mať 384 až 400 strán.\\
our book \textsc{aux} have 384 \textsc{až} 400 pages\\
\glt `Our book will be 384 to 400 pages long.'\\\hfill (Slovak; Andrea Kellö Žačoková, Jana Bašnáková, p.c.)
\z

\ea\label{ex:tom:pl400}
\gll Nasza książka będzie mieć od 384 do 400 stron.\\
our book \textsc{aux} have from 384 to 400 pages\\
\glt `Our book will be 384 to 400 pages long.'
\z


\noindent The differences in the requirements on the scalar dimension with \textit{aż}/\textit{až} in Polish, Czech and Slovak predict that in different languages we may find exclusives specified for different scales. With temporal scales either an inherent ordering is lexicalized (Czech and Slovak \textit{až}) or the context specifies the direction of the scale (so that the prejacent is either later or earlier than expected, as with Polish \textit{dopiero}). And while for %the 
Polish \textit{dopiero} I opted to assume an extra presupposition to account for temporality, for the Slovak \textit{až} I would consider a meaning without a scalar presupposition but encoding scalarity in the assertion (see \citealt{charnavel15} for French scalar opposites of \textit{even} and \textit{only}).

In South Slavic languages, the adverbial particle \textit{čak} has the same semantics as \textit{aż} (\citealt{tomaszewicz13,tomaszewicz13a}), and in Bulgarian, which lacks a counterpart of \textit{dopiero}, \textit{čak} fulfills this role. Just like in Polish, in \REF{ex:tom:95} \textit{čak} contributes the reading that Mary is a noteworthy person for me to talk to. In \REF{ex:tom:96} \textit{čak} would be translated as \textit{dopiero} rather than \textit{aż} to yield the meaning that ‘now’ is a noteworthy time for you to tell me, because it is late, because as we saw in \sectref{sec:tom:4.3} it is more specific: It calls for a contextually relevant ranking of times.

\ea\label{ex:tom:95}
\gll Govorix čak s Mary.\\            
talked \textsc{čak} with Mary\\
\glt `I talked to somebody so important as Mary.’\\\hfill (Bulgarian; \cite[(4)]{tomaszewicz13}, \cite[(1)]{tomaszewicz13a})
\z 

\ea\label{ex:tom:96}
\gll Zašto mi kazvaš čak sega? \\
why me tell \textsc{čak} now\\
\ea `Why are you telling me only now?’
\ex \textit{Unavailable:} `Why are you telling me even now?’\\\hfill (Bulgarian; \cite[(10)]{tomaszewicz13})
\z 
\z 

\noindent Notably, Bosnian-Croatian-Serbian (BCS) lacks a single lexical item with the meaning of \textit{even} (recall the semantics of \textit{even} in \REF{ex:tom:79}, contrasting it with \textit{aż} in \REF{ex:tom:78}). Instead, the two meaning components of \textit{even}, scalarity and additivity, are lexicalized separately by \textit{čak} and \textit{i}, where \textit{i} is the conjunction \textit{and} (\citealt{jasinskajatomaszewicz16}). In example \REF{ex:tom:97a} \textit{čak} appears alone, while in \REF{ex:tom:97b} it is accompanied by \textit{i}, which contributes additivity absent in \REF{ex:tom:97a}, as evidenced by the continuation. 

\ea\label{ex:tom:97}
\ea\label{ex:tom:97a} 
\gll Poslao sam  žalbu čak  dekanu, ali nisam nikome   drugome.\\
sent \textsc{aux} complaint \textsc{čak} dean but \textsc{neg.aux} nobody other\\
\glt `I sent my complaint to somebody so important as the dean, but I did not send it to anyone else.’
\ex\label{ex:tom:97b}    
\gll Poslao sam žalbu čak i dekanu, \minsp{\#} ali nisam nikome   drugome.\\
sent \textsc{aux} complaint \textsc{čak} and dean {} but \textsc{neg.aux} nobody other\\
\glt `I sent my complaint even to the dean \textit{[Incongruous:]} but I did not send it to anyone else.’\hfill (BCS; \cite[(27)]{jasinskajatomaszewicz16})
\z 
\z 



\noindent In \REF{ex:tom:98a} \textit{čak} together with negation yields the reading ‘low on the scale of distance’, while in \REF{ex:tom:98b} the presence of the additive particle (its negated variant \textit{ni}) yields the ‘low on the scale of likelihood’ reading, which is exactly the reading with \textit{even} in its %the 
English translation.

\ea\label{ex:tom:98}
\ea\label{ex:tom:98a} 
\gll Ivan nije čak u Salzburgu, sad je u Zagrebu.\\     
Ivan \textsc{neg.aux} \textsc{čak} in Salzburg now is in Zagreb\\
\glt `Ivan is not as far as Salzburg, he is in Zagreb now.’
\ex\label{ex:tom:98b}    
\gll Ivan nije čak ni u \minsp{\{} Zagrebu / \minsp{\#} Salzburgu\}.\\
Ivan \textsc{neg.aux} \textsc{čak} \textsc{neg}.and in {} Zagreb {} {} Salzburg\\
\glt `Ivan is not even in Zagreb / \textit{[Incongruous:]} Salzburg.’\\\hfill(BCS; \cite[(28--29)]{jasinskajatomaszewicz16})
\z 
\z 

\noindent The facts from BCS %Bosnian-Croatian-Serbian 
suggest that the different meaning components that I identified for the four Polish exclusives can be lexicalized separately, which adds to the other known cases of compositional combination of particles from the exclusive and additive family, e.g. \textit{auch nur} in German \citep{guerzoni03} (\textit{auch} means `also', \textit{nur} means `only', \textit{auch nur} means `even'), \textit{mégis} in Hungarian (Katja Jasinskaja, p.c.; \citealt{csirmazslade20}) (\textit{még} means `still', \textit{is} is an additive particle and \textit{mégis}  is a scalar additive particle).


The case of BCS \textit{čak i} shows that the repertoire of exclusive and additive scalar particles uses the same building blocks at the levels of assertion and presupposition.\footnote{The importance of the distinction between the asserted and presupposed content in the meaning of \textit{even} and \textit{only} goes back to \citealt{horn69}.} The presupposed and asserted components may also end up being swapped, as argued in \citet{guerzoni03} for subspecies of \textit{even} and \textit{only}, or the strength of particular components may differ, e.g. additivity may be weaker as proposed by \citet{crnic11,crnic12} for concessive additive particles. Moreover, the differences in presuppositions have grammatical effects, i.e., affect the distribution of different particles in different syntactic environments \citep{giannakidou07}.

\section{Conclusion}
\label{sec:tom:7}

My analysis of the four Polish exclusives \textit{tylko}, \textit{(za)ledwie}, \textit{dopiero} and \textit{aż} adds to the body of work arguing that we can go beyond the identification of the different flavors of exclusive and additive meanings (as in \citealt{koenig91}) but track the origin of those flavors in their minimal building blocks (\citealt{rullmann97}, \citealt{guerzoni03}, \citealt{lahiri10}, \citealt{crnic11}, \citealt{coppockbeaver11, coppockbeaver10, coppockbeaver14}, a.o.). I identified which meaning components are core and thus reoccur in the different exclusive particles and which components require specific tweaking of parameters. Different settings for specific parameters may influence the flavor of exclusivity, e.g., force the temporal interpretation with \textit{dopiero}, but may also extend their uses beyond that of variants of \textit{only}: I showed that in Polish the restrictions on the alternative set with \textit{(za)ledwie} yield approximative/counterfactual uses, and with \textit{aż} they result in the meaning of \textit{until}.  This work could also be a starting point for the exploration of other expressions with exclusive meaning, such as gradable adjectives, non-focus associating adverbs, etc., for the identification of parameters other than the semantic ingredients that I isolated here. \citet{coppockbeaver11,coppockbeaver10,coppockbeaver14} argue that English adverbial \textit{only} and \textit{just}, and adjectival exclusives \textit{mere}, \textit{sole}, \textit{only}, \textit{single}, \textit{exclusive} vary along two parameters: the ontological type of their arguments and constraints on the question under discussion. \citet{greenberg16,greenberg22} suggests that the variability in the dimensions of the scale could be captured if the scalar presupposition is defined in terms of degrees, with respect to a contextually supplied gradable property. For Polish the next step would be to investigate the building blocks and find a common core with \textit{even}-like particles (\textit{nawet}, which we saw in this chapter, and \textit{choć(by)}), additives (\textit{i}, \textit{też}, \textit{także}, \textit{również}, …), approximatives (the approximative \textit{ledwie} that we saw here belongs to a rich repertoire: \textit{prawie}, \textit{niemal}, \textit{niemalże}, \textit{nieomal}, \textit{omal}, \textit{mało co},…) and exceptives (\textit{oprócz}, \textit{poza}, \textit{ponadto}, \textit{zresztą}, …).





\section*{Abbreviations}

\begin{tabularx}{.5\textwidth}{@{}lQ}
\textsc{aux}&auxiliary\\
\textsc{neg}&negation\\
&\\
\end{tabularx}%
\begin{tabularx}{.5\textwidth}{lQ@{}}
\textsc{top}&topic particle\\
\textsc{Q}&question particle\\
&\\ % this dummy row achieves correct vertical alignment of both tables
\end{tabularx}

\section*{Acknowledgements}
I would like to express my very special thanks to Berit Gehrke and Radek Šimík, without whose assistance and encouragement this chapter would not have been possible. I would also like to thank
the anonymous reviewers for this volume for their insightful comments. The data from languages other than Polish comes from my prior work and its analysis benefited greatly from my exchanges with Katja Jasinskaja, Roumi Pancheva, Anastasija Smirnova, Andrea Kellö Žačoková, Jana Bašnáková, Věra Dvořák, Mojmír Dočekal, Petr Biskup, Denisa Lenertová, Dorothee Fehrmann, Hana Strachoňová, Natasha Fitz\-gibbons, Ivan Kapitonov, Natalia Korotkova, Tania Ionin, Agnieszka Łazorczyk, Agata Renans, Elena Guerzoni, Jim Higginbotham, Sabine Iatridou, Jessica Rett, Barry Schein and Yael Sharvit. A preliminary version of this work was presented at the \textit{Association with focus} workshop at the GLOW 35 conference (University of Potsdam, 31 March 2012) and was greatly enriched by discussions with Roumi Pancheva and by valuable comments from the audience. Any mistakes in this work remain my own.

%\section*{Contributions}
%John Doe contributed to conceptualization, methodology, and validation. 
%Jane Doe contributed to writing of the original draft, review, and editing.

\section*{Literary sources}

Karel Čapek: \textit{Od člověka k člověku III. Čapkova společenská publicistika: Články, glosy, sloupky, úvahy a studie z let 1932--1938. Praha: Československý spisovatel, 1991.}

\sloppy
\printbibliography[heading=subbibliography,notkeyword=this]

\end{document}
